
\setcounter{chapter}{11}

\chapter{主理想整环上的模}

本章的主要目的是证明特别好的环——主理想整环——上有限生成模的结构定理。这个定理是环的理想结构（对主理想整环而言特别简单）反映在其模的结构中的一个例子。如果我们在主理想整环是整数环 \(\mathbb{Z}\) 的情形应用这个结果，就得到有限生成阿贝尔群基本定理的一个证明（我们在第 5 章中考察了它但未加证明）。如果我们代之以主理想整环是系数在域 \(F\) 中的 \(x\) 的多项式环 \(F[x]\) 的情形，就得到矩阵的所谓有理标准形与若尔当标准形的基本结果。在着手证明之前，我们先简要讨论这两个重要应用。

我们已在第 5 章讨论过任何有限生成阿贝尔群都同构于循环阿贝尔群（\(\mathbb{Z}\) 或某个正整数 \(n \neq 0\) 的 \(\mathbb{Z}/n\mathbb{Z}\)）的直和这一结果。还要回忆阿贝尔群与 \(\mathbb{Z}\)-模是一回事。由于 \(\mathbb{Z}\) 的理想正是平凡理想 \((0)\) 与由正整数 \(n\) 生成的主理想 \((n) = n\mathbb{Z}\)，我们看到用模的语言表述的有限生成阿贝尔群基本定理是说：任何有限生成的 \(\mathbb{Z}\)-模都是形如 \(\mathbb{Z}/I\)（\(I\) 是 \(\mathbb{Z}\) 的理想）的模（这些就是循环 \(\mathbb{Z}\)-模）的直和，再加上当直和写成某种特定形式时的唯一性陈述。注意 \(\mathbb{Z}\) 的理想结构与其（有限生成的）模——有限生成阿贝尔群——的结构之间的对应。

有限生成主理想整环上模的基本定理断言：当把主理想整环 \(\mathbb{Z}\) 换成任何主理想整环 \(R\) 时同样的结果成立。特别地，我们在第 10 章已看到，系数在域 \(F\) 中的 \(x\) 的多项式环 \(F[x]\) 上的模与连同固定线性变换 \(T\) 的向量空间 \(V\)（其中元素 \(x\) 作为线性变换 \(T\) 作用在 \(V\) 上）是一回事。这个情形下的基本定理将说：这样的向量空间是形如 \(F[x]/I\)（\(I\) 是 \(F[x]\) 的理想，从而是平凡理想 \((0)\) 或由某个非零多项式 \(f(x)\) 生成的主理想 \((f(x))\)）的模的直和（这些是循环 \(F[x]\)-模），同样再加上当直和写成某种特定形式时的唯一性陈述。若把它翻译回向量空间与线性变换的语言，我们就可以得到关于线性变换 \(T\) 的信息。

例如，假设 \(V\) 是 \(F\) 上 \(n\) 维向量空间，我们为 \(V\) 选取一组基。则给出 \(V\) 到自身的线性变换 \(T\) 与给出系数在 \(F\) 中的 \(n \times n\) 矩阵 \(A\) 是一回事（而为 \(V\) 选取不同的基给出 \(T\) 的另一个矩阵 \(B\)，它与 \(A\) 相似，即形如 \(P^{-1}AP\)，其中 \(P\) 是某个定义基变换的可逆矩阵）。我们将看到，这个情形下的基本定理蕴涵（在域 \(F\) 包含给定线性变换 \(T\) 的所有“特征值”这一假设下）存在 \(V\) 的一组基使得 \(T\) 关联的矩阵尽可能接近对角矩阵，因而具有特别简单的形式。这就是\textbf{若尔当标准形}。\textbf{有理标准形}是 \(T\) 的矩阵的另一种简单形式（它不要求 \(T\) 的特征值属于 \(F\)）。这样我们就能够给出域 \(F\) 上任意 \(n \times n\) 矩阵的标准形，即找出与给定 \(n \times n\) 矩阵相似且特别简单（例如几乎对角）的矩阵。

\begin{example}
设 \(V = \mathbb{Q}^3 = \{(x, y, z) \mid x, y, z \in \mathbb{Q}\}\) 是元素取自有理数域 \(F = \mathbb{Q}\) 的通常的 3 维有序三元组空间，并设 \(T\) 是线性变换
\[
T(x, y, z) = (9x+4y+5z,\ -4x-3z,\ -6x-4y-2z), \quad x, y, z \in \mathbb{Q}.
\]
若取 \(V\) 的标准基 \(e_1 = (1, 0, 0)\)、\(e_2 = (0, 1, 0)\)、\(e_3 = (0, 0, 1)\)，则表示这个线性变换的矩阵 \(A\) 是
\[
A = \begin{pmatrix} 9 & 4 & 5 \\ -4 & 0 & -3 \\ -6 & -4 & -2 \end{pmatrix}.
\]
我们将看到这个矩阵 \(A\) 的若尔当标准形是简单得多的矩阵
\[
B = \begin{pmatrix} 2 & 1 & 0 \\ 0 & 2 & 0 \\ 0 & 0 & 3 \end{pmatrix},
\]
它是通过改取 \(V\) 的基 \(f_1 = (2, -1, -2)\)、\(f_2 = (1, 0, -1)\)、\(f_3 = (3, -2, -2)\) 得到的，因为此时
\[
T(f_1) = T(2, -1, -2) = (4, -2, -4) = 2 \cdot f_1+0 \cdot f_2+0 \cdot f_3,
\]
\[
T(f_2) = T(1, 0, -1) = (4, -1, -4) = 1 \cdot f_1+2 \cdot f_2+0 \cdot f_3,
\]
\[
T(f_3) = T(3, -2, -2) = (9, -6, -6) = 0 \cdot f_1+0 \cdot f_2+3 \cdot f_3,
\]
故关于这组基表示 \(T\) 的矩阵的各列是 \((2, 0, 0)\)、\((1, 2, 0)\) 与 \((0, 0, 3)\)，即 \(T\) 关于这组基的矩阵是 \(B\). 特别地 \(A\) 相似于更简单的矩阵 \(B\).

事实上这个线性变换 \(T\) 不能对角化（即不存在使相应矩阵为对角矩阵的 \(V\) 的基），故矩阵 \(B\) 是 \(T\) 能得到的尽可能接近对角矩阵的矩阵。
\end{example}

下面第一节给出一些一般定义，并在任意主理想整环上陈述和证明基本定理，之后我们回到它关于标准形的应用（关于阿贝尔群的应用见第 5 章）。这些应用可以独立于一般证明来阅读。习题中还概述了另一个在欧几里得整环（从而特别是对环 \(\mathbb{Z}\) 与 \(F[x]\)）上有效、沿行与列运算思路进行、并且有计算价值的证明。

\section{基本理论}

我们首先描述一些一般的有限性条件。设 \(R\) 是环，\(M\) 是左 \(R\)-模。

\begin{definition}
（1）左 \(R\)-模 \(M\) 称为\textbf{诺特 \(R\)-模}，或说它满足\textbf{子模的升链条件}（或子模的 A.C.C.），若不存在无限递增的子模链，即只要
\[
M_1 \subseteq M_2 \subseteq M_3 \subseteq \cdots
\]
是 \(M\) 的子模的递增链，就存在正整数 \(m\) 使得对所有 \(k \geq m\) 有 \(M_k = M_m\)（故这个链在第 \(m\) 步后稳定：\(M_m = M_{m+1} = M_{m+2} = \cdots\)）。

（2）环 \(R\) 称为\textbf{诺特环}，若它作为其自身的左模是诺特的，即 \(R\) 中不存在无限递增的左理想链。
\end{definition}

对（可能非交换的）环 \(R\) 中的右理想与双边理想也可以表述类似的 A.C.C. 概念。对非交换环，这些性质之间不必有联系。

\begin{theorem}\label{thm:12.1}
设 \(R\) 是环，\(M\) 是左 \(R\)-模。则下列命题等价：

（1）\(M\) 是诺特 \(R\)-模。

（2）\(M\) 的每个非空子模集合在包含关系下含有极大元。

（3）\(M\) 的每个子模都是有限生成的。
\end{theorem}

\begin{proof}
〔（1）蕴涵（2）〕假设 \(M\) 是诺特的，\(\Sigma\) 是 \(M\) 的任意非空子模族。任取 \(M_1 \in \Sigma\). 若 \(M_1\) 是 \(\Sigma\) 的极大元，则（2）成立；故假设 \(M_1\) 不是极大的。则存在某个 \(M_2 \in \Sigma\) 使 \(M_1 \subsetneq M_2\). 若 \(M_2\) 是 \(\Sigma\) 的极大元，则（2）成立；故我们可以假设存在某个 \(M_3 \in \Sigma\) 真包含 \(M_2\). 这样继续下去，可以看出若（2）不成立，我们就能用选择公理作出 \(\Sigma\) 中元素的无限严格递增链，与（1）矛盾。

〔（2）蕴涵（3）〕假设（2）成立，\(N\) 是 \(M\) 的任意子模。设 \(\Sigma\) 是 \(N\) 的所有有限生成子模的集合。由于 \(\{0\} \in \Sigma\)，这个集合非空。由（2），\(\Sigma\) 含有极大元 \(N'\). 若 \(N' \neq N\)，取 \(x \in N-N'\). 由于 \(N' \in \Sigma\)，按假设 \(N'\) 是有限生成的，故由 \(N'\) 与 \(x\) 生成的子模也是有限生成的。这与 \(N'\) 的极大性矛盾，故 \(N = N'\) 是有限生成的。

〔（3）蕴涵（1）〕假设（3）成立，\(M_1 \subseteq M_2 \subseteq M_3 \subseteq \cdots\) 是 \(M\) 的子模链。令
\[
N = \bigcup_{i=1}^{\infty} M_i,
\]
并注意 \(N\) 是子模。由（3），\(N\) 是有限生成的，比如说由 \(a_1, a_2, \ldots, a_n\) 生成。由于对所有 \(i\) 有 \(a_i \in N\)，每个 \(a_i\) 都落在链中某个子模中，比如 \(M_{j_i}\). 令 \(m = \max\{j_1, j_2, \ldots, j_n\}\). 则对所有 \(i\) 有 \(a_i \in M_m\)，故它们生成的模包含在 \(M_m\) 中，即 \(N \subseteq M_m\). 这蕴涵对所有 \(k \geq m\) 有 \(M_m = N = M_k\)，这证明了（1）。
\end{proof}

\begin{corollary}\label{cor:12.2}
若 \(R\) 是主理想整环，则 \(R\) 的任何非空理想集合都含有极大元，且 \(R\) 是诺特环。
\end{corollary}

\begin{proof}
主理想整环 \(R\) 在定理中取 \(M = R\) 时满足条件（3）。
\end{proof}

回忆即使 \(M\) 本身是有限生成的 \(R\)-模，\(M\) 的子模也不必是有限生成的，故 \(M\) 是诺特 \(R\)-模这一条件一般比 \(M\) 是有限生成 \(R\)-模这一条件更强。

在转向本章的主要结果之前，我们需要一个关于“线性相关”的结果。

\begin{proposition}\label{pro:12.3}
设 \(R\) 是整环，\(M\) 是秩 \(n < \infty\) 的自由 \(R\)-模。则 \(M\) 的任何 \(n+1\) 个元素都 \(R\)-线性相关，即对任意 \(y_1, y_2, \ldots, y_{n+1} \in M\) 存在不全为零的元素 \(r_1, r_2, \ldots, r_{n+1} \in R\) 使
\[
r_1y_1+r_2y_2+\cdots+r_{n+1}y_{n+1} = 0.
\]
\end{proposition}

\begin{proof}
证明这个结论最快的方法是把 \(R\) 嵌入它的商域 \(F\)（因为 \(R\) 是整环），并注意由于 \(M \cong R \oplus R \oplus \cdots \oplus R\)（\(n\) 次），我们得到 \(M \cong F \oplus F \oplus \cdots \oplus F\)（\(n\) 次）。后者是 \(F\) 上 \(n\) 维向量空间，故 \(M\) 的任何 \(n+1\) 个元素都 \(F\)-线性相关。把纯量的分母清掉（例如两边乘以所有分母的乘积），我们就得到 \(M\) 的这 \(n+1\) 个元素之间的一个 \(R\)-线性相关关系。

另一种做法：设 \(e_1, \ldots, e_n\) 是自由 \(R\)-模 \(M\) 的一组基，\(y_1, \ldots, y_{n+1}\) 是 \(M\) 的任意 \(n+1\) 个元素。对 \(1 \leq i \leq n+1\)，用基 \(e_1, e_2, \ldots, e_n\) 写出 \(y_i = a_{1i}e_1+a_{2i}e_2+\cdots+a_{ni}e_n\). 设 \(A\) 是第 \(i, j\) 个元素为 \(a_{ij}\)（\(1 \leq i \leq n\)，\(1 \leq j \leq n+1\)）且最后一行全为零的 \((n+1) \times (n+1)\) 矩阵，故当然 \(\det A = 0\). 由于 \(R\) 是整环，第 11.4 节推论 \ref{cor:11.27} 表明 \(A\) 的各列 \(R\)-线性相关。\(A\) 的各列的任何相关关系都给出 \(y_i\) 之间的相关关系，完成证明。
\end{proof}

若 \(R\) 是任意整环、\(M\) 是任意 \(R\)-模，回忆
\[
\operatorname{Tor}(M) = \{x \in M \mid \text{对某个非零 } r \in R \text{ 有 } rx = 0\}
\]
是 \(M\) 的子模（称为 \(M\) 的\textbf{挠子模}），而若 \(N\) 是 \(\operatorname{Tor}(M)\) 的任意子模，则 \(N\) 称为 \(M\) 的\textbf{挠子模}（故 \(M\) 的挠子模是所有挠子模的并，即 \(M\) 的极大挠子模）。若 \(\operatorname{Tor}(M) = 0\)，则称模 \(M\) 是\textbf{无挠的}。

对 \(M\) 的任意子模 \(N\)，\(N\) 的\textbf{零化子}是由
\[
\operatorname{Ann}(N) = \{r \in R \mid rn = 0 \text{ 对所有 } n \in N\}
\]
定义的 \(R\) 的理想。注意若 \(N\) 不是 \(M\) 的挠子模，则 \(\operatorname{Ann}(N) = (0)\). 容易看出若 \(N, L\) 是 \(M\) 的子模且 \(N \supseteq L\)，则 \(\operatorname{Ann}(L) \supseteq \operatorname{Ann}(N)\). 若 \(R\) 是主理想整环且 \(N \supseteq L \supseteq M\) 满足 \(\operatorname{Ann}(N) = (a)\)、\(\operatorname{Ann}(L) = (b)\)，则 \(a \mid b\). 特别地，\(M\) 的任何元素 \(x\) 的零化子整除 \(M\) 的零化子（当 \(R = \mathbb{Z}\) 时这由拉格朗日定理推出）。

\begin{definition}
对任意整环 \(R\)，\(R\)-模 \(M\) 的\textbf{秩}是 \(M\) 中 \(R\)-线性无关元素的最大个数。
\end{definition}

上一命题表明对整环上的自由 \(R\)-模 \(M\)，其子模的秩不超过 \(M\) 的秩。这个秩的概念与以前同一术语的用法一致。若环 \(R = F\) 是域，则 \(R\)-模 \(M\) 的秩就是 \(M\) 作为 \(F\) 上向量空间的维数，而任何极大的 \(F\)-线性无关元素组都是 \(M\) 的基。然而对一般的整环，秩为 \(n\) 的 \(R\)-模 \(M\) 不必有“基”，即即使 \(M\) 无挠也不必是自由 \(R\)-模，故对秩的概念必须小心，特别是涉及 \(M\) 的挠元素时。习题 1 到 6 和习题 20 给出秩的另一种刻画，并给出一些（秩为 1 的）无挠但非自由的 \(R\)-模的例子。

下一个重要结果表明，若 \(N\) 是主理想整环上有限秩自由模的子模，则 \(N\) 仍是有限秩自由模，而且可以为这两个模选取以简单方式相互联系的生成元。

\begin{theorem}\label{thm:12.4}
设 \(R\) 是主理想整环，\(M\) 是秩 \(n\) 的自由 \(R\)-模，\(N\) 是 \(M\) 的子模。则

（1）\(N\) 是秩 \(m \leq n\) 的自由模，且

（2）存在 \(M\) 的一组基 \(y_1, y_2, \ldots, y_n\)，使得 \(a_1y_1, a_2y_2, \ldots, a_my_m\) 是 \(N\) 的一组基，其中 \(a_1, a_2, \ldots, a_m\) 是 \(R\) 的非零元素，且满足整除关系 \(a_1 \mid a_2 \mid \cdots \mid a_m\)。
\end{theorem}

\begin{proof}
对 \(N = \{0\}\) 定理是平凡的，故假设 \(N \neq \{0\}\). 对 \(M\) 到 \(R\) 的每个 \(R\)-模同态 \(\varphi\)，\(N\) 的像 \(\varphi(N)\) 是 \(R\) 的子模，即 \(R\) 的理想。由于 \(R\) 是主理想整环，这个理想必是主的，比如说 \(\varphi(N) = (a_\varphi)\)（某个 \(a_\varphi \in R\)）。设
\[
\Sigma = \{(a_\varphi) \mid \varphi \in \operatorname{Hom}_R(M, R)\}
\]
为这样由 \(M\) 到 \(R\) 的 \(R\)-模同态得到的 \(R\) 中主理想的集合。这个集合当然非空，因为取 \(\varphi\) 为平凡同态就表明 \((0) \in \Sigma\). 由推论 \ref{cor:12.2}，\(\Sigma\) 至少有一个极大元，即至少存在一个 \(M\) 到 \(R\) 的同态 \(\nu\)，使得主理想 \(\nu(N) = (a_\nu)\) 不真包含于 \(\Sigma\) 的任何其他元素中。对这个极大元令 \(a_1 = a_\nu\)，并设 \(y \in N\) 是在同态 \(\nu\) 下映到生成元 \(a_1\) 的元素：\(\nu(y) = a_1\).

现在我们证明元素 \(a_1\) 非零。设 \(x_1, x_2, \ldots, x_n\) 是自由模 \(M\) 的任意一组基，\(\pi_i \in \operatorname{Hom}_R(M, R)\) 是到关于这组基的第 \(i\) 个坐标的自然投影同态。由于 \(N \neq \{0\}\)，存在某个 \(i\) 使 \(\pi_i(N) \neq 0\)，这特别表明 \(\Sigma\) 包含的不只是平凡理想 \((0)\). 由于 \((a_1)\) 是 \(\Sigma\) 的极大元，可知 \(a_1 \neq 0\).

接下来我们证明这个元素 \(a_1\) 整除每个 \(\varphi \in \operatorname{Hom}_R(M, R)\) 的 \(\varphi(y)\). 为此设 \(d\) 是由 \(a_1\) 与 \(\varphi(y)\) 生成的主理想的生成元。则 \(d\) 是 \(R\) 中 \(a_1\) 与 \(\varphi(y)\) 的公因子，且对某些 \(r_1, r_2 \in R\) 有 \(d = r_1a_1+r_2\varphi(y)\). 考虑同态 \(\psi = r_1\nu+r_2\varphi\)（从 \(M\) 到 \(R\)）。则 \(\psi(y) = (r_1\nu+r_2\varphi)(y) = r_1a_1+r_2\varphi(y) = d\)，故 \(d \in \psi(N)\)，从而也有 \((d) \subseteq \psi(N)\). 但 \(d\) 是 \(a_1\) 的因子，故也有 \((a_1) \subseteq (d)\). 于是 \((a_1) \subseteq (d) \subseteq \psi(N)\)，而由 \((a_1)\) 的极大性必有等号：\((a_1) = (d) = \psi(N)\). 特别地 \((a_1) = (d)\) 表明 \(a_1 \mid \varphi(y)\).

若把这一点应用于投影同态 \(\pi_i\)，我们看到 \(a_1\) 整除所有 \(\pi_i(y)\). 对某些 \(b_i \in R\) 写 \(\pi_i(y) = a_1b_i\)（\(1 \leq i \leq n\)），并定义
\[
y_1 = \sum_{i=1}^n b_ix_i.
\]
注意 \(a_1y_1 = y\). 由于 \(a_1 = \nu(y) = \nu(a_1y_1) = a_1\nu(y_1)\) 且 \(a_1\) 是整环 \(R\) 的非零元素，这表明
\[
\nu(y_1) = 1.
\]
现在我们验证这个元素 \(y_1\) 可以作为 \(M\) 的一组基中的一个元素、而 \(a_1y_1\) 可以作为 \(N\) 的一组基中的一个元素，即：

（a）\(M = Ry_1 \oplus \ker\nu\)，且

（b）\(N = Ra_1y_1 \oplus (N \cap \ker\nu)\).

为看出（a），设 \(x\) 是 \(M\) 中任意元素，并写 \(x = \nu(x)y_1+(x-\nu(x)y_1)\). 由于
\[
\nu(x-\nu(x)y_1) = \nu(x)-\nu(x)\nu(y_1) = \nu(x)-\nu(x) \cdot 1 = 0,
\]
我们看到 \(x-\nu(x)y_1\) 是 \(\nu\) 的核中的元素。这表明 \(x\) 可以写成 \(Ry_1\) 中元素与 \(\nu\) 的核中元素之和，故 \(M = Ry_1+\ker\nu\). 为看出这个和是直和，假设 \(ry_1\) 也在 \(\nu\) 的核中。则 \(0 = \nu(ry_1) = r\nu(y_1) = r\) 表明这个元素确实是 0.

对于（b），注意由 \(a_1\) 作为 \(\nu(N)\) 的生成元的定义，对每个 \(x' \in N\) 都有 \(a_1 \mid \nu(x')\). 若对某个 \(b \in R\) 写 \(\nu(x') = ba_1\)，则我们在（a）中用过的分解给出 \(x' = \nu(x')y_1+(x'-\nu(x')y_1) = ba_1y_1+(x'-ba_1y_1)\)，其中第二个被加项在 \(\nu\) 的核中且是 \(N\) 的元素。这表明 \(N = Ra_1y_1+(N \cap \ker\nu)\). （b）中和的直接性由（a）中和的直接性推出。

现在我们通过对 \(N\) 的秩 \(m\) 归纳证明定理的第（1）部分。若 \(m = 0\)，则 \(N\) 是挠模，从而 \(N = 0\)（因为自由模无挠），故（1）平凡地成立。于是假设 \(m > 0\). 由于上面（b）中的和是直和，容易看出 \(N \cap \ker\nu\) 的秩为 \(m-1\)（参见习题 3）。由归纳假设，\(N \cap \ker\nu\) 是秩 \(m-1\) 的自由 \(R\)-模。再次由（b）中和的直接性可知把 \(a_1y_1\) 加到 \(N \cap \ker\nu\) 的任意基上就得到 \(N\) 的基，故 \(N\) 也是自由的（秩为 \(m\)），这证明了（1）。

最后我们通过对 \(M\) 的秩 \(n\) 归纳证明（2）。对子模 \(\ker\nu\) 应用（1）可知这个子模是自由的，而由（a）中和的直接性它是秩 \(n-1\) 的。对模 \(\ker\nu\)（它扮演 \(M\) 的角色）及其子模 \(\ker\nu \cap N\)（它扮演 \(N\) 的角色）应用归纳假设，我们看到存在 \(\ker\nu\) 的一组基 \(y_2, y_3, \ldots, y_n\)，使得 \(a_2y_2, a_3y_3, \ldots, a_my_m\) 是 \(N \cap \ker\nu\) 的基，其中 \(a_2, a_3, \ldots, a_m\) 是 \(R\) 的元素且 \(a_2 \mid a_3 \mid \cdots \mid a_m\). 由于（a）与（b）中的和都是直和，\(y_1, y_2, \ldots, y_n\) 是 \(M\) 的基，而 \(a_1y_1, a_2y_2, \ldots, a_my_m\) 是 \(N\) 的基。为完成归纳，剩下要证明 \(a_1\) 整除 \(a_2\). 通过在 \(M\) 的基上定义 \(\varphi(y_1) = \varphi(y_2) = 1\)、对所有 \(i > 2\) 定义 \(\varphi(y_i) = 0\)，定义从 \(M\) 到 \(R\) 的同态 \(\varphi\). 则对这个同态 \(\varphi\) 有 \(a_1 = \varphi(a_1y_1)\)，故 \(a_1 \in \varphi(N)\)，从而 \((a_1) \subseteq \varphi(N)\). 由 \((a_1)\) 在 \(\Sigma\) 中的极大性可知 \((a_1) = \varphi(N)\). 由于 \(a_2 = \varphi(a_2y_2) \in \varphi(N)\)，我们就有 \(a_2 \in (a_1)\)，即 \(a_1 \mid a_2\). 这完成了定理的证明。
\end{proof}

回忆左 \(R\)-模 \(C\) 是\textbf{循环 \(R\)-模}（对任意环 \(R\)，不必交换也不必带 1），若存在元素 \(x \in C\) 使 \(C = Rx\). 我们于是可以定义 \(R\)-模同态
\[
\pi : R \longrightarrow C, \qquad \pi(r) = rx,
\]
它由假设 \(C = Rx\) 是满射。第一同构定理给出（左）\(R\)-模的同构
\[
R/\ker\pi \cong C.
\]
若 \(R\) 是主理想整环，则 \(\ker\pi\) 是主理想 \((a)\)，故我们看到循环 \(R\)-模 \(C\) 具有 \(R/(a)\) 的形式，其中 \((a) = \operatorname{Ann}(C)\).

循环模是最简单的模（因为它们只需要一个生成元）。基本定理的存在性部分断言主理想整环上任何有限生成模都同构于有限多个循环模的直和。

\begin{theorem}\label{thm:12.5}
（\textbf{基本定理，存在性：不变因子形式}）设 \(R\) 是主理想整环，\(M\) 是有限生成的 \(R\)-模。

（1）则 \(M\) 同构于有限多个循环模的直和。更精确地说，
\[
M \cong R^r \oplus R/(a_1) \oplus R/(a_2) \oplus \cdots \oplus R/(a_m)
\]
对某个整数 \(r \geq 0\) 与 \(R\) 的某些非零元素 \(a_1, a_2, \ldots, a_m\)，它们不是 \(R\) 中的单位，且满足整除关系
\[
a_1 \mid a_2 \mid \cdots \mid a_m.
\]

（2）\(M\) 无挠当且仅当 \(M\) 是自由模。

（3）在（1）的分解中，\(\operatorname{Tor}(M) \cong R/(a_1) \oplus \cdots \oplus R/(a_m)\). 特别地 \(M\) 是挠模当且仅当 \(r = 0\)，此时 \(M\) 的零化子是理想 \((a_m)\)。
\end{theorem}

\begin{proof}
按假设 \(M\) 可以由一个有限元素集生成，故设 \(x_1, x_2, \ldots, x_n\) 是 \(M\) 的基数最小的生成元组。设 \(R^n\) 是以 \(b_1, b_2, \ldots, b_n\) 为基的秩 \(n\) 自由 \(R\)-模，并通过定义 \(\pi(b_i) = x_i\)（对所有 \(i\)）定义同态 \(\pi : R^n \to M\)，由于 \(x_1, \ldots, x_n\) 生成 \(M\)，它自动是满射。由模的第一同构定理有 \(R^n/\ker\pi \cong M\). 现在对 \(R^n\) 及其子模 \(\ker\pi\) 应用定理 \ref{thm:12.4}，我们可以选取 \(R^n\) 的另一组基 \(y_1, y_2, \ldots, y_n\)，使得 \(a_1y_1, a_2y_2, \ldots, a_my_m\) 是 \(\ker\pi\) 的基，其中 \(a_1, a_2, \ldots, a_m\) 是 \(R\) 的元素且 \(a_1 \mid a_2 \mid \cdots \mid a_m\). 这蕴涵
\[
M \cong R^n/\ker\pi \cong (Ry_1 \oplus \cdots \oplus Ry_n)/(Ra_1y_1 \oplus \cdots \oplus Ra_my_m).
\]
为辨认右端的商，我们用把 \((a_1y_1, \ldots, a_ny_n)\) 映到 \((a_1 \bmod (a_1), \ldots, a_m \bmod (a_m), a_{m+1}, \ldots, a_n)\) 的自然满 \(R\)-模同态
\[
Ry_1 \oplus \cdots \oplus Ry_n \longrightarrow R/(a_1) \oplus \cdots \oplus R/(a_m) \oplus R \oplus \cdots \oplus R.
\]
这个映射的核显然是满足 \(a_i\) 整除 \(a_i\)（\(i = 1, 2, \ldots, m\)）的那些元素组成的集合，即 \(Ra_1y_1 \oplus Ra_2y_2 \oplus \cdots \oplus Ra_my_m\)（参见习题 7）。于是我们得到
\[
M \cong R/(a_1) \oplus \cdots \oplus R/(a_m) \oplus R^{n-m}.
\]
若 \(a\) 是 \(R\) 中的单位，则 \(R/(a) = 0\)，故在这个直和中我们可以去掉开头任何为单位的 \(a_i\). 这给出（1）中的分解（取 \(r = n-m\)）。

由于对 \(R\) 的任何非零元素 \(a\)，\(R/(a)\) 是挠 \(R\)-模，（1）立即蕴涵 \(M\) 无挠当且仅当 \(M \cong R^r\)，这就是（2）。

第（3）部分由定义立即得到，因为 \(R/(a)\) 的零化子显然是理想 \((a)\)。
\end{proof}

我们很快将证明定理 \ref{thm:12.5} 中分解的唯一性，即若我们有
\[
M \cong R^{r'} \oplus R/(b_1) \oplus R/(b_2) \oplus \cdots \oplus R/(b_{m'})
\]
（某个整数 \(r' \geq 0\) 与 \(R\) 的某些非零且非单位的元素 \(b_1, b_2, \ldots, b_{m'}\)，满足 \(b_1 \mid b_2 \mid \cdots \mid b_{m'}\)），则 \(r = r'\)、\(m = m'\) 且对所有 \(i\) 有 \((a_i) = (b_i)\)（故 \(a_i = b_i\) 在单位意义下）。正是整除条件 \(a_1 \mid a_2 \mid \cdots \mid a_m\) 给出了这个唯一性。

\begin{definition}
定理 \ref{thm:12.5} 中的整数 \(r\) 称为 \(M\) 的\textbf{自由秩}或\textbf{Betti 数}，而 \(R\) 中的元素 \(a_1, a_2, \ldots, a_m\)（在乘以 \(R\) 中单位的意义下定义）称为 \(M\) 的\textbf{不变因子}。
\end{definition}

注意在证明 \(M\) 的不变因子唯一之前，我们应当恰当地说 \(M\) 的\textbf{一组}不变因子（自由秩也是如此）——我们的意思是指给出上面定理（1）中那样的 \(M\) 的分解的任何元素。

利用中国剩余定理可以把定理 \ref{thm:12.5} 中的循环模进一步分解，使 \(M\) 是其零化子尽可能简单（即 \((0)\) 或由 \(R\) 中素元的幂生成）的循环模的直和。这给出另一种分解，我们也将看到它是唯一的，现在描述如下。

假设 \(a\) 是主理想整环 \(R\) 的非零元素。则由于 \(R\) 是唯一分解整环，我们可以写
\[
a = up_1^{\alpha_1}p_2^{\alpha_2}\cdots p_s^{\alpha_s},
\]
其中 \(p_i\) 是 \(R\) 中互不相同的素元，\(u\) 是单位。这个分解在单位意义下唯一，故理想 \((p_i^{\alpha_i})\)（\(i = 1, \ldots, s\)）有唯一的定义。对 \(i \neq j\) 我们有 \((p_i^{\alpha_i})+(p_j^{\alpha_j}) = R\)，因为这两个理想的和由一个最大公因子生成，而对互不相同的素元 \(p_i, p_j\) 这个最大公因子是 1. 换句话说，理想 \((p_i^{\alpha_i})\)（\(i = 1, \ldots, s\)）两两互余。所有这一切理想的交是理想 \((a)\)，因为 \(a\) 是 \(p_1^{\alpha_1}, p_2^{\alpha_2}, \ldots, p_s^{\alpha_s}\) 的最小公倍。则中国剩余定理（定理 7.17）表明作为环也作为 \(R\)-模有
\[
R/(a) \cong R/(p_1^{\alpha_1}) \oplus R/(p_2^{\alpha_2}) \oplus \cdots \oplus R/(p_s^{\alpha_s}).
\]
把它应用于定理 \ref{thm:12.5} 中的模，使我们能把每个直和项 \(R/(a_i)\)（\(a_i\) 是 \(M\) 的不变因子）写成零化子为 \(a_i\) 的素幂因子的循环模的直和。这证明了：

\begin{theorem}\label{thm:12.6}
（\textbf{基本定理，存在性：初等因子形式}）设 \(R\) 是主理想整环，\(M\) 是有限生成的 \(R\)-模。则 \(M\) 是有限多个零化子为 \((0)\) 或由 \(R\) 中素元的幂生成的循环模的直和，即
\[
M \cong R^r \oplus R/(p_1^{\alpha_1}) \oplus R/(p_2^{\alpha_2}) \oplus \cdots \oplus R/(p_s^{\alpha_s}),
\]
其中 \(r \geq 0\) 是整数，\(p_1^{\alpha_1}, \ldots, p_s^{\alpha_s}\) 是 \(R\) 中（不必互异的）素元的正幂。
\end{theorem}

我们通过 \(M\) 的不变因子的素幂因子证明了定理 \ref{thm:12.6}。事实上我们将看到，把 \(M\) 分解为如定理 \ref{thm:12.6} 中那样零化子为 \((0)\) 或素幂的循环模的直和是唯一的，即整数 \(r\) 与理想 \((p_1^{\alpha_1}), \ldots, (p_s^{\alpha_s})\) 对 \(M\) 唯一确定。这些素幂有一个名称：

\begin{definition}
设 \(R\) 是主理想整环，\(M\) 是如定理 \ref{thm:12.6} 中的有限生成 \(R\)-模。素幂 \(p_1^{\alpha_1}, \ldots, p_s^{\alpha_s}\)（在乘以 \(R\) 中单位的意义下定义）称为 \(M\) 的\textbf{初等因子}。
\end{definition}

假设 \(M\) 是主理想整环 \(R\) 上有限生成的挠模。若对定理 \ref{thm:12.6} 的分解中出现的互异素元 \(p_1, p_2, \ldots, p_n\)，我们把对应于同一个 \(p_i\) 的所有循环因子归在一起，就特别看到 \(M\) 可以写成直和
\[
M = N_1 \oplus N_2 \oplus \cdots \oplus N_n,
\]
其中 \(N_i\) 由 \(M\) 中所有被素元 \(p_i\) 的某个幂零化的元素组成。这个结果对 \(R\) 上可能非有限生成的模也成立：

\begin{theorem}\label{thm:12.7}
（\textbf{准素分解定理}）设 \(R\) 是主理想整环，\(M\) 是非零挠 \(R\)-模（不必有限生成），其零化子 \(a\) 非零。假设 \(a\) 在 \(R\) 中分解为互异素幂之积为
\[
a = up_1^{\alpha_1}p_2^{\alpha_2}\cdots p_n^{\alpha_n},
\]
并设 \(N_i = \{x \in M \mid p_i^{\alpha_i}x = 0\}\)（\(1 \leq i \leq n\)）。则 \(N_i\) 是 \(M\) 的零化子为 \(p_i^{\alpha_i}\) 的子模，且是 \(M\) 中所有被 \(p_i\) 的某个幂零化的元素组成的子模。我们有
\[
M = N_1 \oplus N_2 \oplus \cdots \oplus N_n.
\]
若 \(M\) 是有限生成的，则每个 \(N_i\) 是有限多个零化子为 \(p_i^{\alpha_i}\) 的因子的循环模的直和。
\end{theorem}

\begin{proof}
在 \(M\) 于 \(R\) 上有限生成的情形我们已经证明了这些结果。在一般情形，显然 \(N_i\) 是 \(M\) 的零化子整除 \(p_i^{\alpha_i}\) 的子模。由于 \(R\) 是主理想整环，对 \(i \neq j\) 理想 \((p_i^{\alpha_i})\) 与 \((p_j^{\alpha_j})\) 互余，故通过修改中国剩余定理证明中的论证并把它应用于模，就容易证明 \(M\) 的直和分解。用这个直和分解容易看出 \(N_i\) 的零化子恰为 \(p_i^{\alpha_i}\)。
\end{proof}

\begin{definition}
上一定理中的子模 \(N_i\) 称为 \(M\) 的 \textbf{\(p_i\)-准素分支}。
\end{definition}

注意用这个术语，有限生成模 \(M\) 的初等因子正是 \(\operatorname{Tor}(M)\) 的各准素分支的不变因子。

现在我们证明关于基本定理中分解的唯一性陈述。

注意若 \(M\) 是交换环 \(R\) 上的任意模、\(a\) 是 \(R\) 的元素，则 \(aM = \{am \mid m \in M\}\) 是 \(M\) 的子模。还要回忆在主理想整环 \(R\) 中非零素理想是极大的，故 \(R\) 关于非零素理想的商是域。

\begin{lemma}\label{lem:12.8}
设 \(R\) 是主理想整环，\(p\) 是 \(R\) 中的素元。用 \(F\) 表示域 \(R/(p)\)。

（1）设 \(M = R^r\). 则 \(M/pM \cong F^r\).

（2）设 \(M = R/(a)\)，其中 \(a\) 是 \(R\) 的非零元素。则 \(M/pM \cong F\)（若 \(p\) 整除 \(a\)），\(M/pM = 0\)（若 \(p\) 不整除 \(a\)）。

（3）设 \(M = R/(a_1) \oplus R/(a_2) \oplus \cdots \oplus R/(a_k)\)，其中每个 \(a_i\) 都被 \(p\) 整除。则 \(M/pM \cong F^k\).
\end{lemma}

\begin{proof}
（1）存在从 \(R^r\) 到 \((R/(p))^r\) 的自然映射，由把 \((a_1, \ldots, a_r)\) 映到 \((a_1 \bmod (p), \ldots, a_r \bmod (p))\) 给出。显然这是满 \(R\)-模同态，其核由所有坐标都被 \(p\) 整除的 \(r\) 元组组成，即 \(pR^r\)，故 \(R^r/pR^r \cong (R/(p))^r\)，这就是（1）。

（2）这由同构定理推出：首先注意 \(p(R/(a))\) 是理想 \((p)\) 在商 \(R/(a)\) 中的像，从而是 \(((p)+(a))/(a)\). 理想 \((p)+(a)\) 由 \(p\) 与 \(a\) 的最大公因子生成，从而若 \(p\) 整除 \(a\) 则是 \((p)\)，否则是 \(R\). 于是若 \(p\) 整除 \(a\) 则 \(pM = (p)/(a)\)，否则 \(pM = R/(a) = M\).

若 \(p\) 整除 \(a\)，则 \(M/pM = (R/(a))/((p)/(a)) \cong R/(p)\)；若 \(p\) 不整除 \(a\)，则 \(M/pM = M/M = 0\)，这证明了（2）。

（3）这由（2）按定理 \ref{thm:12.5} 第（1）部分的证明方式推出。
\end{proof}

\begin{theorem}\label{thm:12.9}
（\textbf{基本定理，唯一性}）设 \(R\) 是主理想整环。

（1）两个有限生成 \(R\)-模 \(M_1\) 与 \(M_2\) 同构当且仅当它们有相同的自由秩和相同的不变因子表。

（2）两个有限生成 \(R\)-模 \(M_1\) 与 \(M_2\) 同构当且仅当它们有相同的自由秩和相同的初等因子表。
\end{theorem}

\begin{proof}
若 \(M_1\) 与 \(M_2\) 有相同的自由秩以及相同的不变因子表或相同的初等因子表，则它们显然同构。

假设 \(M_1\) 与 \(M_2\) 同构。\(M_1\) 与 \(M_2\) 之间的任何同构都把 \(M_1\) 中的挠元素映到 \(M_2\) 中的挠元素，故必有 \(\operatorname{Tor}(M_1) \cong \operatorname{Tor}(M_2)\). 于是 \(R^{r_1} \cong M_1/\operatorname{Tor}(M_1) \cong M_2/\operatorname{Tor}(M_2) \cong R^{r_2}\)，其中 \(r_1\) 是 \(M_1\) 的自由秩、\(r_2\) 是 \(M_2\) 的自由秩。设 \(p\) 是 \(R\) 中任意非零素元。则由 \(R^{r_1} \cong R^{r_2}\) 得到 \(R^{r_1}/pR^{r_1} \cong R^{r_2}/pR^{r_2}\). 由上一引理的（1），这蕴涵 \(F^{r_1} \cong F^{r_2}\)，其中 \(F\) 是域 \(R/pR\). 于是我们有 \(F\) 上 \(r_1\) 维向量空间与 \(F\) 上 \(r_2\) 维向量空间的同构，故 \(r_1 = r_2\)，\(M_1\) 与 \(M_2\) 有相同的自由秩。

我们化归为证明 \(M_1\) 与 \(M_2\) 有相同的不变因子表与初等因子表。为此只需处理同构的挠模 \(\operatorname{Tor}(M_1)\) 与 \(\operatorname{Tor}(M_2)\)，即我们可以不失一般性地假设 \(M_1\) 与 \(M_2\) 都是挠 \(R\)-模。

我们先证明它们有相同的初等因子。只需证明对任意固定的素元 \(p\)，凡是 \(p\) 的幂的初等因子对 \(M_1\) 与 \(M_2\) 都相同。若 \(M_1 \cong M_2\)，则 \(M_1\) 的 \(p\)-准素子模（\(=\) 初等因子为 \(p\) 的幂的那些循环因子的直和）同构于 \(M_2\) 的 \(p\)-准素子模，因为这些是被 \(p\) 的某个幂零化的元素组成的子模。因此我们化归为证明：若两个零化子为 \(p\) 的幂的模 \(M_1\) 与 \(M_2\) 同构，则它们有相同的初等因子。

我们对 \(M_1\) 的零化子中 \(p\) 的幂归纳（由于 \(M_1\) 与 \(M_2\) 同构，这与 \(M_2\) 的零化子相同）。若这个幂为 0，则 \(M_1\) 与 \(M_2\) 都是 0，结论成立。否则 \(M_1\)（与 \(M_2\)）有非平凡的初等因子。假设 \(M_1\) 的初等因子为
\[
\underbrace{p, p, \ldots, p}_{m \text{ 个}}, p^{\alpha_1}, p^{\alpha_2}, \ldots, p^{\alpha_s},
\]
其中 \(2 \leq \alpha_1 \leq \alpha_2 \leq \cdots \leq \alpha_s\)，即 \(M_1\) 是循环模的直和，其生成元设为 \(x_1, x_2, \ldots, x_m, x_{m+1}, \ldots, x_{m+s}\)，零化子分别为 \((p), (p), \ldots, (p), (p^{\alpha_1}), \ldots, (p^{\alpha_s})\). 则子模 \(pM_1\) 的初等因子为 \(p^{\alpha_1-1}, p^{\alpha_2-1}, \ldots, p^{\alpha_s-1}\)，因为 \(pM_1\) 是循环模的直和，其生成元为 \(px_1, px_2, \ldots, px_m, px_{m+1}, \ldots, px_{m+s}\)，零化子分别为 \((1), (1), \ldots, (1), (p^{\alpha_1-1}), \ldots, (p^{\alpha_s-1})\). 类似地，若 \(M_2\) 的初等因子为
\[
\underbrace{p, p, \ldots, p}_{n \text{ 个}}, p^{\beta_1}, p^{\beta_2}, \ldots, p^{\beta_t},
\]
其中 \(2 \leq \beta_1 \leq \beta_2 \leq \cdots \leq \beta_t\)，则 \(pM_2\) 的初等因子为 \(p^{\beta_1-1}, p^{\beta_2-1}, \ldots, p^{\beta_t-1}\).

由于 \(M_1 \cong M_2\)，也有 \(pM_1 \cong pM_2\)，而 \(pM_1\) 的零化子中 \(p\) 的幂比 \(M_1\) 的零化子中 \(p\) 的幂少 1. 由归纳假设，\(pM_1\) 的初等因子与 \(pM_2\) 的初等因子相同，即 \(s = t\) 且 \(\alpha_i = \beta_i\)（\(i = 1, 2, \ldots, s\)）。由上一引理的（3），我们有 \(m+s = n+t\)，从而由已看到的 \(s = t\) 得到 \(m = n\). 这证明 \(M_1\) 的初等因子集合与 \(M_2\) 的初等因子集合相同。

现在我们证明 \(M_1\) 与 \(M_2\) 必有相同的不变因子。假设 \(a_1 \mid a_2 \mid \cdots \mid a_m\) 是 \(M_1\) 的不变因子。通过取这些元素的素幂因子，我们得到 \(M_1\) 的一组初等因子。注意此时不变因子上的整除关系蕴涵：\(a_m\) 是这些初等因子中各素元最大幂的乘积，\(a_{m-1}\) 是去掉 \(a_m\) 的因子后剩余初等因子中各素元最大幂的乘积，依此类推。若 \(b_1 \mid b_2 \mid \cdots \mid b_n\) 是 \(M_2\) 的不变因子，则同样通过取这些元素的素幂因子得到 \(M_2\) 的一组初等因子。但上面我们证明了 \(M_1\) 与 \(M_2\) 的初等因子相同，由此可知不变因子也相同。
\end{proof}

\begin{corollary}\label{cor:12.10}
设 \(R\) 是主理想整环，\(M\) 是有限生成的 \(R\)-模。

（1）\(M\) 的初等因子是 \(M\) 的不变因子的素幂因子。

（2）\(M\) 的最大不变因子是 \(M\) 的初等因子中各互异素元最大幂的乘积，次大不变因子是剩余初等因子中各互异素元最大幂的乘积，依此类推。
\end{corollary}

\begin{proof}
（1）中的做法给出一组初等因子，而由定理，\(M\) 的初等因子唯一，故（1）中的做法给出初等因子集合。（2）类似。
\end{proof}

\begin{corollary}\label{cor:12.11}
（\textbf{有限生成阿贝尔群基本定理}）见定理 5.3 与定理 5.5。
\end{corollary}

\begin{proof}
在定理 \ref{thm:12.5}、\ref{thm:12.6} 与 \ref{thm:12.9} 中取 \(R = \mathbb{Z}\)（但注意为计算方便，第 5 章中不变因子是按相反次序列的）。
\end{proof}

推论 \ref{cor:12.10} 中在初等因子与不变因子之间转换的步骤，在有限生成阿贝尔群的情形下第 5 章已有相当详细的描述。

还要注意若把有限生成模 \(M\) 写成形如 \(R/(a)\) 的循环模的直和，则除非加上某些附加条件（例如不变因子的整除条件，或 \(a\) 是素元的幂这一条件），出现的理想 \((a)\) 一般不是唯一的。为判断两个模是否同构，必须先把它们写成这样的标准（或典范）形式。

\subsection*{习题}

\begin{enumerate}
\item 设 \(M\) 是整环 \(R\) 上的模。

（a）假设 \(x\) 是 \(M\) 中的非零挠元素。证明 \(x\) 与 0 “线性相关”。由此断定 \(\operatorname{Tor}(M)\) 的秩为 0，从而特别地任何挠 \(R\)-模的秩为 0.

（b）证明 \(M\) 的秩等于（无挠的）商 \(M/\operatorname{Tor}(M)\) 的秩。

\item 设 \(M\) 是整环 \(R\) 上的模。

（a）假设 \(M\) 的秩为 \(n\)，\(x_1, x_2, \ldots, x_n\) 是 \(M\) 的任意一组极大线性无关元素。设 \(N = Rx_1+\cdots+Rx_n\) 是它们生成的子模。证明 \(N \cong R^n\)，且商 \(M/N\) 是挠 \(R\)-模（等价地：\(x_1, \ldots, x_n\) 线性无关，且对任意 \(y \in M\) 存在非零元素 \(r \in R\) 使 \(ry\) 可以写成 \(x_i\) 的线性组合 \(r_1x_1+\cdots+r_nx_n\)）。

（b）反之证明：若 \(M\) 含有一个秩 \(n\) 的自由子模 \(N\)（即 \(N \cong R^n\)）使得商 \(M/N\) 是挠 \(R\)-模，则 \(M\) 的秩为 \(n\).〔设 \(y_1, y_2, \ldots, y_{n+1}\) 是 \(M\) 的任意 \(n+1\) 个元素。利用 \(M/N\) 是挠模这一事实，对某些非零元素 \(r_1, \ldots, r_{n+1} \in R\) 把 \(r_iy_i\) 写成 \(N\) 的一组基的线性组合。用命题 \ref{pro:12.3} 证明中的论证看出 \(r_iy_i\) 从而 \(y_i\) 线性相关。〕

\item 设 \(R\) 是整环，\(A\) 与 \(B\) 是秩分别为 \(m\) 与 \(n\) 的 \(R\)-模。证明 \(A \oplus B\) 的秩为 \(m+n\).〔利用上一习题。〕

\item 设 \(R\) 是整环，\(M\) 是 \(R\)-模，\(N\) 是 \(M\) 的子模。假设 \(M\) 的秩为 \(n\)、\(N\) 的秩为 \(r\)、商 \(M/N\) 的秩为 \(s\). 证明 \(n = r+s\).〔设 \(x_1, x_2, \ldots, x_s\) 是 \(M\) 中在 \(M/N\) 中的像是极大线性无关组的元素，\(x_{s+1}, x_{s+2}, \ldots, x_{s+r}\) 是 \(N\) 中的极大线性无关组。证明 \(x_1, x_2, \ldots, x_{s+r}\) 在 \(M\) 中线性无关，且对任意 \(y \in M\) 存在非零 \(r \in R\) 使 \(ry\) 是这些元素的线性组合。然后用习题 2。〕

\item 设 \(R = \mathbb{Z}[x]\)，\(M = (2, x)\) 是 \(2\) 与 \(x\) 生成的理想，看作 \(R\) 的子模。证明 \(\{2, x\}\) 不是 \(M\) 的基。〔找出这两个元素之间的一个非平凡 \(R\)-线性相关关系。〕证明 \(M\) 的秩为 1，但 \(M\) 不是秩 1 的自由模（参见习题 2）。

\item 证明若 \(R\) 是整环且 \(M\) 是 \(R\) 的任意非主理想，则 \(M\) 是无挠的、秩为 1，但不是自由 \(R\)-模。

\item 设 \(R\) 是任意环，\(A_1, A_2, \ldots, A_m\) 是 \(R\)-模，\(B_i\) 是 \(A_i\) 的子模（\(1 \leq i \leq m\)）。证明
\[
(A_1 \oplus A_2 \oplus \cdots \oplus A_m)/(B_1 \oplus B_2 \oplus \cdots \oplus B_m) \cong (A_1/B_1) \oplus (A_2/B_2) \oplus \cdots \oplus (A_m/B_m).
\]

\item 设 \(R\) 是主理想整环，\(B\) 是挠 \(R\)-模，\(p\) 是 \(R\) 中的素元。证明若对某个非零 \(b \in B\) 有 \(pb = 0\)，则 \(\operatorname{Ann}(B) \subseteq (p)\).

\item 给出一个整环 \(R\) 与非零挠 \(R\)-模 \(M\) 的例子，使 \(\operatorname{Ann}(M) = 0\). 证明若 \(N\) 是有限生成的挠 \(R\)-模，则 \(\operatorname{Ann}(N) \neq 0\).

\item 对主理想整环 \(R\) 中的素元 \(p\) 与 \(R\)-模 \(N\)，证明 \(N\) 的 \(p\)-准素分支是 \(N\) 的子模，并证明 \(N\) 是它的所有 \(p\)-准素分支的直和（这样的分支不必只有有限多个）。

\item 设 \(R\) 是主理想整环，\(a\) 是 \(R\) 的非零元素，\(M = R/(a)\). 证明对 \(R\) 的任意素元 \(p\)，
\[
p^kM/p^{k+1}M \cong \begin{cases} F, & \text{若 } k < n, \\ 0, & \text{若 } k \geq n, \end{cases}
\]
其中 \(n\) 是 \(p\) 在 \(R\) 中整除 \(a\) 的幂，\(F = R/(p)\).

\item 设 \(R\) 是主理想整环，\(p\) 是 \(R\) 中的素元。

（a）设 \(M\) 是有限生成的挠 \(R\)-模。用上一习题证明 \(p^{k-1}M/p^kM \cong F^{n_k}\)，其中 \(F\) 是域 \(R/(p)\)，而 \(n_k\) 是 \(M\) 的形如 \(p^a\)（\(a \geq k\)）的初等因子的个数。

（b）假设 \(M_1\) 与 \(M_2\) 是同构的有限生成的挠 \(R\)-模。用（a）证明对每个 \(k \geq 0\)，\(M_1\) 与 \(M_2\) 有相同个数的形如 \(p^a\)（\(a \geq k\)）的初等因子。证明这蕴涵 \(M_1\) 与 \(M_2\) 有相同的初等因子集合。

\item 若 \(M\) 是主理想整环 \(R\) 上有限生成的模，描述 \(M/\operatorname{Tor}(M)\) 的结构。

\item 设 \(R\) 是主理想整环，\(M\) 是挠 \(R\)-模。证明 \(M\) 不可约（参见第 10.3 节习题 9 到 11）当且仅当对任意非零元素 \(m \in M\) 有 \(M = Rm\)，且 \(m\) 的零化子是非零素理想 \((p)\).

\item 证明若 \(R\) 是诺特环，则 \(R^n\) 是诺特 \(R\)-模。〔固定 \(R^n\) 的一组基。若 \(M\) 是 \(R^n\) 的子模，证明 \(M\) 中元素的第一坐标的集合是 \(R\) 的理想，从而是有限生成的。设 \(m_1, m_2, \ldots, m_k\) 是 \(M\) 中第一坐标生成这个理想的元素。证明 \(M\) 的任意元素都可以写成 \(m_1, \ldots, m_k\) 的 \(R\)-线性组合加上 \(M \cap R^{n-1}\) 中的元素之和，其中 \(R^{n-1}\) 是 \(R^n\) 中第一坐标为零的元素组成的子模，然后用对 \(n\) 的归纳。〕
\end{enumerate}

\noindent 下面一组习题用涉及行与列运算的矩阵论证概述定理 \ref{thm:12.5} 在 \(R\) 是欧几里得整环这一特殊情形下的一个证明。这特别适用于应用中有意义的 \(R = \mathbb{Z}\) 与 \(R = F[x]\) 的情形，并且有计算价值。

\noindent 设 \(R\) 是欧几里得整环，\(M\) 是 \(R\)-模。

\begin{enumerate}
\setcounter{enumi}{15}
\item 证明 \(M\) 有限生成当且仅当对某个整数 \(n\) 存在满 \(R\)-同态 \(\varphi : R^n \to M\)（这对任何环 \(R\) 都成立）。

\noindent 假设 \(\varphi : R^n \to M\) 是满 \(R\)-模同态。由习题 15，\(\ker\varphi\) 是有限生成的。若 \(x_1, x_2, \ldots, x_n\) 是 \(R^n\) 的一组基，\(y_1, \ldots, y_m\) 是 \(\ker\varphi\) 的生成元，我们有
\[
y_i = a_{i1}x_1+a_{i2}x_2+\cdots+a_{in}x_n, \quad i = 1, 2, \ldots, m,
\]
其中系数 \(a_{ij} \in R\). 由此同态 \(\varphi\)（从而 \(M\) 的模结构）由 \(R^n\) 的生成元的选取与矩阵 \(A = (a_{ij})\) 的选取确定。这样的矩阵 \(A\) 称为\textbf{关系矩阵}。

\item （a）证明在 \(R^n\) 的基中交换 \(x_i\) 与 \(x_j\)，相当于在相应关系矩阵中交换第 \(i\) 列与第 \(j\) 列。

（b）证明对任意 \(a \in R\)，在 \(R^n\) 的基中把元素 \(x_j\) 换成 \(x_j-ax_i\) 得到 \(R^n\) 的另一组基，且关于这组基的相应关系矩阵与原关系矩阵相同，只是把第 \(j\) 列的 \(a\) 倍加到了第 \(i\) 列上。〔注意 \(\cdots+a_ix_i+\cdots+a_jx_j+\cdots = \cdots+(a_i+aa_j)x_i+\cdots+a_j(x_j-ax_i)+\cdots\)。〕

\item （a）证明交换生成元 \(y_i\) 与 \(y_j\) 相当于在关系矩阵中交换第 \(i\) 行与第 \(j\) 行。

（b）证明对任意 \(a \in R\)，把元素 \(y_j\) 换成 \(y_j-ay_i\) 得到 \(\ker\varphi\) 的另一组生成元，且关于这组生成元的相应关系矩阵与原关系矩阵相同，只是把第 \(i\) 行的 \(-a\) 倍加到了第 \(j\) 行上。

\item 由前两个习题，通过为 \(R^n\) 与 \(\ker\varphi\) 选取不同的生成元，我们可以对给定的关系矩阵施行初等行运算与列运算。若所有关系矩阵都是零矩阵，则 \(\ker\varphi = 0\) 且 \(M \cong R^n\). 否则设 \(a_1\) 是 \(M\) 的一个固定初始关系矩阵中所有元素的（非零）最大公因子（回忆 \(R\) 是欧几里得整环）。

（a）证明通过初等行运算与列运算，我们可以使 \(a_1\) 出现在形如 \(\begin{pmatrix} a_1 & 0 \\ 0 & A' \end{pmatrix}\) 的关系矩阵中，其中 \(a_1\) 整除所有 \(a_{ij}\)（\(i = 1, 2, \ldots, m\)，\(j = 1, 2, \ldots, n\)）。

（b）证明存在形如 \(\begin{pmatrix} a_1 & 0 \\ 0 & a_2 \\ & & A''\end{pmatrix}\) 的关系矩阵，其中 \(a_1\) 整除它的所有元素。

（c）设 \(a_2\) 是（b）中关系矩阵中除 \(a_1\) 之外所有元素的最大公因子。证明存在形如 \(\begin{pmatrix} a_1 & 0 & 0 \\ 0 & a_2 & 0 \\ 0 & 0 & A'''\end{pmatrix}\) 的关系矩阵，其中 \(a_1\) 整除 \(a_2\)，且 \(a_2\) 整除矩阵中其他所有元素。

（d）证明存在形如 \(\begin{pmatrix} D & 0 \\ 0 & 0 \end{pmatrix}\) 的关系矩阵，其中 \(D\) 是主对角线元素为非零元素 \(a_1, a_2, \ldots, a_k\)（\(k \leq n\)）的对角矩阵，满足 \(a_1 \mid a_2 \mid \cdots \mid a_k\). 由此断定
\[
M \cong R/(a_1) \oplus R/(a_2) \oplus \cdots \oplus R/(a_k) \oplus R^{n-k}.
\]

\noindent 若 \(n\) 不是 \(M\) 所需生成元的最小个数，则上面开头若干个元素 \(a_1, a_2, \ldots\) 中会有一些是单位，上述相应的直和项为 0. 若去掉这些无关的因子，我们就得到了模 \(M\) 的不变因子。进一步，对应于上述直和项的 \(R^n\) 的新生成元的像将是 \(M\) 在其不变因子分解中各个循环子模的一组 \(R\)-生成元（注意对应于 \(a_i\) 为单位的因子的生成元在 \(M\) 中的像是 0）。关系矩阵约化中施行的列运算对应于如习题 17 所描述的、改变用于 \(R^n\) 的基的运算：

（a）交换第 \(i\) 列与第 \(j\) 列对应于交换 \(R^n\) 的基中第 \(i\) 个与第 \(j\) 个元素。

（b）对任意 \(a \in R\)，把第 \(j\) 列的 \(a\) 倍加到第 \(i\) 列上对应于从第 \(j\) 个基元素中减去第 \(i\) 个基元素的 \(a\) 倍。

\noindent 因此，记录所施行的列运算并同样地改变 \(M\) 的初始生成元的选取，就给出一组 \(M\) 在其不变因子分解中各个循环子模的 \(R\)-生成元。

\noindent 一旦确定了 \(M\) 的初始生成元组与初始关系矩阵，这个过程在计算上是相当快的。元素 \(a_1\) 用欧几里得算法作为初始关系矩阵中元素的最大公因子来确定。使用行与列运算，我们可以得到这些元素的适当的线性组合，使这个最大公因子出现在新关系矩阵的 \((1,1)\) 位置。然后减去第一列与第一行的适当倍数，得到如习题 19（b）中的矩阵，再迭代这个过程。下一节末尾给出了这个步骤在某些特殊情形下的例子。

\item 设 \(R\) 是商域为 \(F\) 的整环，\(M\) 是任意 \(R\)-模。证明 \(M\) 的秩等于 \(F\) 上向量空间 \(F \otimes_R M\) 的维数。

\item 证明主理想整环上有限生成的模是投射模当且仅当它是自由模。

\item 设 \(R\) 是不是域的主理想整环。证明没有有限生成的 \(R\)-模是内射模。〔用第 10.5 节习题 4 分别考虑挠模与自由模。〕
\end{enumerate}

\section{有理标准形}

现在我们把这些关于有限生成模的结果应用于主理想整环是系数在域 \(F\) 中的 \(x\) 的多项式环 \(F[x]\) 这一特殊情形。

设 \(V\) 是 \(F\) 上 \(n\) 维有限维向量空间，\(T\) 是 \(V\) 的固定线性变换（即从 \(V\) 到自身）。正如我们在第 10 章看到的，我们可以把 \(V\) 看作 \(F[x]\)-模，其中元素 \(x\) 作为线性变换 \(T\) 作用在 \(V\) 上（从而 \(x\) 的任何多项式都作为相应 \(T\) 的多项式作用在 \(V\) 上）。由于按假设 \(V\) 在 \(F\) 上有限维，按定义它是有限生成的 \(F\)-模，从而当然是有限生成的 \(F[x]\)-模，故上一节的分类定理适用。

任何非零的自由 \(F[x]\)-模（同构于若干份 \(F[x]\) 的直和）都是 \(F\) 上的无限维向量空间，故若 \(V\) 在 \(F\) 上有限维，它就事实上必是挠 \(F[x]\)-模（即其自由秩为 0）。由基本定理可知，此时 \(V\) 作为 \(F[x]\)-模同构于循环挠 \(F[x]\)-模的直和。我们将看到 \(V\) 的这个分解使我们能为 \(V\) 选取一组基，使得线性变换 \(T\) 关于它的矩阵表示具有特定的简单形式。当我们使用 \(V\) 的不变因子分解时，就得到 \(T\) 的矩阵的\textbf{有理标准形}，本节我们分析它。当我们使用初等因子分解（且 \(F\) 包含 \(T\) 的所有特征值）时，就得到若尔当标准形，它在下一节考虑，并在前面提到过是表示 \(T\) 的、尽可能接近对角矩阵的矩阵。基本定理的唯一性部分保证有理标准形与若尔当标准形是唯一的（这正是它们被称为标准形的原因）。

这些标准形的一个重要用途是分类 \(V\) 的不同的线性变换。特别地，它们使我们能判断两个矩阵何时表示同一个线性变换，即两个给定的 \(n \times n\) 矩阵何时相似。

注意这又是一个用被作用的空间的结构（例如 \(V\) 的不变因子分解）来获得关于起作用的代数对象（这里是线性变换）的重要信息的例子。第 18 章将考察群作用在向量空间上的情形（它以“群的表示论”之名著称）。

在详细描述有理标准形之前，我们首先引入一些线性代数内容。

\begin{definition}
（1）\(F\) 的元素 \(\lambda\) 称为线性变换 \(T\) 的\textbf{特征值}，若存在非零向量 \(v \in V\) 使 \(T(v) = \lambda v\). 此时称 \(v\) 是 \(T\) 的、对应于特征值 \(\lambda\) 的\textbf{特征向量}。

（2）若 \(A\) 是系数在 \(F\) 中的 \(n \times n\) 矩阵，则元素 \(\lambda\) 称为 \(A\) 的、对应特征向量 \(v\) 的\textbf{特征值}，若 \(v\) 是非零 \(n \times 1\) 列向量且 \(Av = \lambda v\).

（3）若 \(\lambda\) 是线性变换 \(T\) 的特征值，则集合 \(\{v \in V \mid T(v) = \lambda v\}\) 称为 \(T\) 对应于特征值 \(\lambda\) 的\textbf{特征空间}。类似地，若 \(\lambda\) 是 \(n \times n\) 矩阵 \(A\) 的特征值，则满足 \(Av = \lambda v\) 的 \(n \times 1\) 矩阵 \(v\) 的集合称为 \(A\) 对应于特征值 \(\lambda\) 的特征空间。
\end{definition}

注意若固定 \(V\) 的一组基 \(\mathcal{B}\)，则 \(V\) 的任何线性变换 \(T\) 都有一个关联的 \(n \times n\) 矩阵 \(A\). 反之，若 \(A\) 是任意 \(n \times n\) 矩阵，则由对 \(v \in V\) 定义 \(T(v) = Av\)（右端的 \(v\) 是由 \(v\) 关于 \(V\) 的固定基 \(\mathcal{B}\) 的坐标组成的 \(n \times 1\) 向量）定义的映射 \(T\) 是 \(V\) 的线性变换。则 \(v\) 是 \(T\) 的、特征值为 \(\lambda\) 的特征向量当且仅当 \(v\) 关于 \(\mathcal{B}\) 的坐标向量是 \(A\) 的、特征值为 \(\lambda\) 的特征向量。换句话说，线性变换 \(T\) 的特征值与 \(T\) 关于 \(V\) 的任何固定基的矩阵 \(A\) 的特征值相同。

\begin{definition}
从 \(V\) 到 \(V\) 的线性变换的\textbf{行列式}是表示这个线性变换的任何矩阵的行列式（注意这不依赖于所用基的选取）。
\end{definition}

\begin{proposition}\label{pro:12.12}
下列命题等价：

（1）\(\lambda\) 是 \(T\) 的特征值；

（2）\(\lambda I-T\) 是 \(V\) 的奇异线性变换；

（3）\(\det(\lambda I-T) = 0\).
\end{proposition}

\begin{proof}
由于 \(\lambda\) 是 \(T\) 的、对应特征向量 \(v\) 的特征值当且仅当 \(v\) 是 \(\lambda I-T\) 的核中的非零向量，可知（1）与（2）等价。（2）与（3）由我们关于行列式的结果等价。
\end{proof}

\begin{definition}
设 \(x\) 是 \(F\) 上的未定元。多项式 \(\det(xI-T)\) 称为 \(T\) 的\textbf{特征多项式}，记作 \(c_T(x)\). 若 \(A\) 是系数在 \(F\) 中的 \(n \times n\) 矩阵，则 \(\det(xI-A)\) 称为 \(A\) 的特征多项式，记作 \(c_A(x)\).
\end{definition}

通过展开行列式容易看出 \(T\) 或 \(A\) 的特征多项式是次数为 \(n = \dim V\) 的首一多项式。命题 \ref{pro:12.12} 说 \(T\)（或 \(A\)）的特征值集合正是 \(T\)（相应地 \(A\)）的特征多项式的根集合。特别地，\(T\) 至多有 \(n\) 个不同的特征值。

我们已经看到，\(V\) 经由线性变换 \(T\) 看作 \(F[x]\)-模时是挠 \(F[x]\)-模。设 \(m(x) \in F[x]\) 是生成 \(V\) 在 \(F[x]\) 中的零化子的唯一首一多项式。等价地，\(m(x)\) 是零化 \(V\) 的（即满足 \(m(T)\) 是零线性变换的）次数最小的唯一首一多项式，且若 \(f(x) \in F[x]\) 是零化 \(V\) 的任何多项式，则 \(m(x)\) 整除 \(f(x)\). 由于 \(F\) 上所有 \(n \times n\) 矩阵的环同构于 \(V\) 到自身的所有线性变换的集合（通过为 \(V\) 选取一组基得到同构），可知对 \(F\) 上任何 \(n \times n\) 矩阵 \(A\)，同样存在唯一的、以 \(m(A)\) 为零矩阵的次数最小的首一多项式。

\begin{definition}
生成 \(F[x]\) 中理想 \(\operatorname{Ann}(V)\) 的唯一首一多项式称为 \(T\) 的\textbf{极小多项式}，记作 \(m_T(x)\). 在矩阵 \(A\) 处取值时给出零矩阵的次数最小的唯一首一多项式称为 \(A\) 的\textbf{极小多项式}，记作 \(m_A(x)\).
\end{definition}

容易看出（参见习题 5）这些极小多项式的次数至多为 \(n^2\)，其中 \(n\) 是 \(V\) 的维数。我们很快就会证明 \(T\) 的极小多项式整除 \(T\) 的特征多项式（这就是凯莱--哈密顿定理），对 \(A\) 类似，故事实上这些多项式的次数至多为 \(n\).

现在我们描述线性变换 \(T\)（相应地 \(n \times n\) 矩阵 \(A\)）的有理标准形。由定理 \ref{thm:12.5} 我们有 \(F[x]\)-模的同构
\begin{equation}
V \cong F[x]/(a_1(x)) \oplus F[x]/(a_2(x)) \oplus \cdots \oplus F[x]/(a_m(x)), \tag{12.1}
\end{equation}
其中 \(a_1(x), a_2(x), \ldots, a_m(x)\) 是 \(F[x]\) 中次数至少为 1 的多项式，且满足整除条件
\[
a_1(x) \mid a_2(x) \mid \cdots \mid a_m(x).
\]
这些不变因子 \(a_i(x)\) 只在 \(F[x]\) 中单位的意义下确定，但由于 \(F[x]\) 的单位正是 \(F\) 的非零元素（即非零常数多项式），我们可以通过要求它们首一来使这些多项式唯一。

由于 \(V\) 的零化子是理想 \((a_m(x))\)（定理 \ref{thm:12.5} 第（3）部分），我们立即得到：

\begin{proposition}\label{pro:12.13}
极小多项式 \(m_T(x)\) 是 \(V\) 的最大不变因子。\(V\) 的所有不变因子都整除 \(m_T(x)\).
\end{proposition}

下面我们将看到如何不仅计算 \(T\) 的极小多项式，而且计算其他不变因子。

现在我们为上面（12.1）分解中 \(V\) 的每个直和项选取一组基，使得 \(T\) 的矩阵相当简单。回忆作用在（12.1）左端的线性变换 \(T\) 就是元素 \(x\) 在（12.1）同构右端每个因子上由乘法作用。

我们在第 11 章命题 1 之后的例子中看到，元素 \(1, \overline{x}, \overline{x^2}, \ldots, \overline{x^{k-1}}\) 给出向量空间 \(F[x]/(a(x))\) 的一组基，其中 \(a(x) = x^k+b_{k-1}x^{k-1}+\cdots+b_1x+b_0\) 是 \(F[x]\) 中任意首一多项式，\(\overline{x} = x \bmod (a(x))\). 关于这组基，乘以 \(x\) 的线性变换作用方式很简单：
\[
\overline{1} \mapsto \overline{x}, \qquad \overline{x} \mapsto \overline{x^2}, \qquad \ldots, \qquad \overline{x^{k-2}} \mapsto \overline{x^{k-1}}, \qquad \overline{x^{k-1}} \mapsto \overline{x^k} = -b_0-b_1\overline{x}-\cdots-b_{k-1}\overline{x^{k-1}},
\]
其中最后一个等式是因为在 \(F[x]/(a(x))\) 中 \(\overline{a(x)} = 0\)，即 \(\overline{x^k}+b_{k-1}\overline{x^{k-1}}+\cdots+b_1\overline{x}+b_0 = 0\). 关于这组基，乘以 \(x\) 的矩阵因此是
\[
\begin{pmatrix}
0 & 0 & \cdots & 0 & -b_0\\
1 & 0 & \cdots & 0 & -b_1\\
0 & 1 & \cdots & 0 & -b_2\\
\vdots & & \ddots & & \vdots\\
0 & 0 & \cdots & 1 & -b_{k-1}
\end{pmatrix}.
\]
这样的矩阵有一个名称：

\begin{definition}
设 \(a(x) = x^k+b_{k-1}x^{k-1}+\cdots+b_1x+b_0\) 是 \(F[x]\) 中的任意首一多项式。\(a(x)\) 的\textbf{伴侣矩阵}是第一条次对角线上全为 1、最后一列自上而下为 \(-b_0, -b_1, \ldots, -b_{k-1}\)、其余位置全为零的 \(k \times k\) 矩阵。\(a(x)\) 的伴侣矩阵记作 \(C_{a(x)}\).
\end{definition}

我们把这一点应用于上面（12.1）右端的每个循环模，并设 \(\mathcal{B}_i\) 是 \(V\) 中对应于在同构（12.1）下为循环因子 \(F[x]/(a_i(x))\) 所选基的元素。则按定义线性变换 \(T\) 通过 \(a_i(x)\) 的伴侣矩阵作用于 \(\mathcal{B}_i\)，因为我们已经看到乘以 \(x\) 就是这样作用的。由于（12.1）右端的和是直和，各 \(\mathcal{B}_i\) 的并 \(\mathcal{B}\) 给出 \(V\) 的一组基，而关于这组基线性变换 \(T\) 的矩阵是不变因子的伴侣矩阵的直和，即
\begin{equation}
\begin{pmatrix} C_{a_1(x)} & & & \\ & C_{a_2(x)} & & \\ & & \ddots & \\ & & & C_{a_m(x)} \end{pmatrix}. \tag{12.2}
\end{equation}
注意这个矩阵由 \(F[x]\)-模 \(V\) 的不变因子唯一确定，而由定理 \ref{thm:12.9}，不变因子表在同构意义下唯一确定模 \(V\)（作为 \(F[x]\)-模）。

\begin{definition}
（1）若一个矩阵是次数至少为 1 的首一多项式 \(a_1(x), \ldots, a_m(x)\) 的伴侣矩阵的直和，且 \(a_1(x) \mid a_2(x) \mid \cdots \mid a_m(x)\)，则称它处于\textbf{有理标准形}。多项式 \(a_i(x)\) 称为这个矩阵的\textbf{不变因子}。这样的矩阵也称为以 \(a_i(x)\) 的伴侣矩阵为块的\textbf{分块对角}矩阵。

（2）线性变换 \(T\) 的\textbf{有理标准形}是表示 \(T\) 的、处于有理标准形的矩阵。
\end{definition}

我们已经看到任何线性变换 \(T\) 都有有理标准形。现在我们看到这个有理标准形是唯一的（因此称为 \(T\) 的有理标准形）。为看出这一点，注意我们从直和分解确定 \(T\) 的矩阵所用的过程是可逆的。假设 \(b_1(x), b_2(x), \ldots, b_l(x)\) 是 \(F[x]\) 中次数至少为 1 的首一多项式，且对所有 \(i\) 有 \(b_i(x) \mid b_{i+1}(x)\)，并假设对 \(V\) 的某组基 \(\mathcal{E}\)，\(T\) 关于 \(\mathcal{E}\) 的矩阵是 \(b_i(x)\) 的伴侣矩阵的直和。则 \(V\) 必是 \(T\)-稳定子空间 \(D_i\) 的直和，每个 \(b_i(x)\) 对应一个，使得 \(T\) 在每个 \(D_i\) 上的矩阵是 \(b_i(x)\) 的伴侣矩阵。设 \(\mathcal{E}_i\) 是 \(D_i\) 相应的（有序）基（故 \(\mathcal{E}\) 是各 \(\mathcal{E}_i\) 的并），\(e_i\) 是 \(\mathcal{E}_i\) 中第一个基元素。则容易看出 \(D_i\) 是以 \(e_i\) 为生成元的循环 \(F[x]\)-模，且 \(D_i\) 的零化子是 \(b_i(x)\). 于是挠 \(F[x]\)-模 \(V\) 以两种方式分解为循环 \(F[x]\)-模的直和，二者都满足定理 \ref{thm:12.5} 的条件，即二者都给出不变因子表。由于由定理 \ref{thm:12.9} 不变因子唯一，\(a_i(x)\) 与 \(b_i(x)\) 必相差 \(F[x]\) 中的一个单位因子，而由假设这些多项式都是首一的，故对所有 \(i\) 必有 \(a_i(x) = b_i(x)\). 这证明了下面结果：

\begin{theorem}\label{thm:12.14}
（\textbf{线性变换的有理标准形}）设 \(V\) 是域 \(F\) 上的有限维向量空间，\(T\) 是 \(V\) 的线性变换。

（1）存在 \(V\) 的一组基，使得 \(T\) 关于它的矩阵处于有理标准形，即它是以次数至少为 1 的首一多项式 \(a_1(x), a_2(x), \ldots, a_m(x)\)（满足 \(a_1(x) \mid a_2(x) \mid \cdots \mid a_m(x)\)）的伴侣矩阵为对角块的分块对角矩阵。

（2）\(T\) 的有理标准形唯一。
\end{theorem}

用“有理”一词是为了表明这个标准形完全在域 \(F\) 内计算，且对任何线性变换 \(T\) 都存在。若尔当标准形（稍后考虑）则不是这样，它只在域 \(F\) 包含 \(T\) 的特征值时才存在（另见推论 \ref{cor:12.18} 之后的注记）。

下面这个结果把相似的线性变换（即相差一个基变换的同一个线性变换）的概念翻译成模的语言，并把这个概念与有理标准形联系起来。

\begin{theorem}\label{thm:12.15}
设 \(S\) 与 \(T\) 是 \(V\) 的线性变换。则下列命题等价：

（1）\(S\) 与 \(T\) 是相似的线性变换；

（2）经由 \(S\) 与经由 \(T\) 从 \(V\) 得到的 \(F[x]\)-模是同构的 \(F[x]\)-模；

（3）\(S\) 与 \(T\) 有相同的有理标准形。
\end{theorem}

\begin{proof}
〔（1）蕴涵（2）〕假设存在非奇异线性变换 \(U\) 使 \(S = UTU^{-1}\). 向量空间同构 \(U : V \to V\) 也是 \(F[x]\)-模同态，其中 \(x\) 在第一个 \(V\) 上经由 \(T\) 作用、在第二个 \(V\) 上经由 \(S\) 作用，因为例如 \(U(xv) = U(Tv) = UT(v) = SU(v) = x(Uv)\). 故这是（2）中两个模的 \(F[x]\)-模同构。

〔（2）蕴涵（3）〕假设（2）成立，并用 \(V_1\) 表示经由 \(S\) 使 \(V\) 成为的 \(F[x]\)-模、用 \(V_2\) 表示经由 \(T\) 使 \(V\) 成为的 \(F[x]\)-模。由于作为 \(F[x]\)-模有 \(V_1 \cong V_2\)，它们有相同的不变因子表。于是 \(S\) 与 \(T\) 有共同的有理标准形。

〔（3）蕴涵（1）〕假设（3）成立。由于按假设 \(S\) 与 \(T\) 关于 \(V\) 的某组（可能不同的）基有相同的矩阵表示，它们相差一个基变换是 \(V\) 的同一个线性变换，从而是相似的。
\end{proof}

设 \(A\) 是元素取自 \(F\) 的任意 \(n \times n\) 矩阵，\(V\) 是 \(F\) 上 \(n\) 维向量空间。回忆我们可以通过为 \(V\) 选取一组基并令 \(T(v) = Av\)（右端的 \(v\) 表示 \(v\) 关于所选基的坐标组成的 \(n \times 1\) 列向量）来定义 \(V\) 上的线性变换 \(T\)（这不过是线性变换与矩阵的通常等同）。则（当然）这个 \(T\) 关于这组基的矩阵就是给定的矩阵 \(A\). 换句话说，元素取自域 \(F\) 的任何 \(n \times n\) 矩阵都作为某个 \(n\) 维向量空间上某个线性变换 \(T\) 的矩阵出现。

向量空间的线性变换与矩阵之间的这个字典使我们能用矩阵的语言陈述前面两个结果：

\begin{theorem}\label{thm:12.16}
（\textbf{矩阵的有理标准形}）设 \(A\) 是域 \(F\) 上的 \(n \times n\) 矩阵。

（1）\(A\) 相似于一个处于有理标准形的矩阵，即存在 \(F\) 上可逆的 \(n \times n\) 矩阵 \(P\) 使 \(P^{-1}AP\) 是以次数至少为 1 的首一多项式 \(a_1(x), a_2(x), \ldots, a_m(x)\)（满足 \(a_1(x) \mid a_2(x) \mid \cdots \mid a_m(x)\)）的伴侣矩阵为对角块的分块对角矩阵。

（2）\(A\) 的有理标准形唯一。
\end{theorem}

\begin{definition}
域 \(F\) 上 \(n \times n\) 矩阵的\textbf{不变因子}是它的有理标准形的不变因子。
\end{definition}

\begin{theorem}\label{thm:12.17}
设 \(A\) 与 \(B\) 是域 \(F\) 上的 \(n \times n\) 矩阵。则 \(A\) 与 \(B\) 相似当且仅当 \(A\) 与 \(B\) 有相同的有理标准形。
\end{theorem}

若 \(A\) 是元素取自域 \(F\) 的矩阵，而 \(F\) 是更大域 \(K\) 的子域，则我们也可以把 \(A\) 看作 \(K\) 上的矩阵。下一个结果表明 \(A\) 的有理标准形与相似性问题不依赖于包含 \(A\) 的元素的域。

\begin{corollary}\label{cor:12.18}
设 \(A\) 与 \(B\) 是域 \(F\) 上的两个 \(n \times n\) 矩阵，并假设 \(F\) 是域 \(K\) 的子域。

（1）\(A\) 的有理标准形在 \(K\) 上计算与在 \(F\) 上计算是相同的。\(A\) 的极小多项式、特征多项式与不变因子在把 \(A\) 看作 \(F\) 上矩阵与看作 \(K\) 上矩阵时都相同。

（2）\(A\) 与 \(B\) 在 \(K\) 上相似当且仅当它们在 \(F\) 上相似，即存在元素取自 \(K\) 的可逆 \(n \times n\) 矩阵 \(P\) 使 \(B = P^{-1}AP\)，当且仅当存在元素取自 \(F\) 的（一般不同的）可逆 \(n \times n\) 矩阵 \(Q\) 使 \(B = Q^{-1}AQ\).
\end{corollary}

\begin{proof}
（1）设 \(M\) 是在较小的域 \(F\) 上计算时 \(A\) 的有理标准形。由于 \(M\) 满足在 \(K\) 上有理标准形定义中的条件，有理标准形的唯一性蕴涵 \(M\) 也是 \(A\) 在 \(K\) 上的有理标准形。于是把 \(A\) 看作 \(F\) 上或 \(K\) 上矩阵时其不变因子都相同。特别地，由于极小多项式是 \(A\) 的最大不变因子，它也不依赖于把 \(A\) 看作哪个域上的矩阵。由特征多项式的行列式定义显然这个多项式只依赖于 \(A\) 的元素（我们很快会看到特征多项式是 \(A\) 的所有不变因子之积，这将给出这个结果的另一个证明）。

（2）若 \(A\) 与 \(B\) 在较小的域 \(F\) 上相似，则它们在 \(K\) 上显然相似。反之，若 \(A\) 与 \(B\) 在 \(K\) 上相似，则它们在 \(K\) 上有相同的有理标准形。由（1），它们在 \(F\) 上有相同的有理标准形，从而由定理 \ref{thm:12.17} 在 \(F\) 上相似。
\end{proof}

这个推论特别断言 \(n \times n\) 矩阵 \(A\) 的有理标准形是元素在包含 \(A\) 的元素的最小域中的 \(n \times n\) 矩阵。进一步，即使我们允许用元素取自更大域的非奇异矩阵来共轭 \(A\)，这个标准形仍是同一个矩阵。这解释了“有理标准形”这一术语。

下一个命题给出矩阵（或线性变换）的特征多项式与其不变因子之间的联系，对确定这些不变因子相当有用（特别是对小尺寸的矩阵）。

\begin{lemma}\label{lem:12.19}
设 \(a(x) \in F[x]\) 是任意首一多项式。

（1）\(a(x)\) 的伴侣矩阵的特征多项式是 \(a(x)\).

（2）若 \(M\) 是由矩阵 \(A_1, A_2, \ldots, A_k\) 的直和给出的分块对角矩阵，则 \(M\) 的特征多项式是 \(A_1, A_2, \ldots, A_k\) 的特征多项式之积。
\end{lemma}

\begin{proof}
这两个都是直接的习题。
\end{proof}

\begin{proposition}\label{pro:12.20}
设 \(A\) 是域 \(F\) 上的 \(n \times n\) 矩阵。

（1）\(A\) 的特征多项式是 \(A\) 的所有不变因子之积。

（2）（\textbf{凯莱--哈密顿定理}）\(A\) 的极小多项式整除 \(A\) 的特征多项式。

（3）\(A\) 的特征多项式整除 \(A\) 的极小多项式的某个幂。特别地，这两个多项式（不计重数）有相同的根。

若把矩阵 \(A\) 换成 \(F\) 上 \(n\) 维向量空间的线性变换 \(T\)，同样的结论成立。
\end{proposition}

\begin{proof}
设 \(B\) 是 \(A\) 的有理标准形。由上一引理，\(B\) 的分块对角形式表明 \(B\) 的特征多项式是 \(A\) 的不变因子的伴侣矩阵的特征多项式之积。由上面引理的第一部分，\(a(x)\) 的伴侣矩阵 \(C_{a(x)}\) 的特征多项式就是 \(a(x)\)，这蕴涵 \(B\) 的特征多项式是 \(A\) 的不变因子之积。由于 \(A\) 与 \(B\) 相似，它们有相同的特征多项式，这证明了（1）。断言（2）由（1）立即推出，因为 \(A\) 的极小多项式是 \(A\) 的最大不变因子。所有不变因子都整除最大不变因子这一事实立即蕴涵（3）。最后一个断言由向量空间的线性变换与矩阵之间的字典显然得出。
\end{proof}

注意命题的第（2）部分就是断言矩阵 \(A\) 满足它自己的特征多项式，即作为矩阵有 \(c_A(A) = 0\)，这是凯莱--哈密顿定理的通常表述。还要注意它蕴涵 \(A\) 的极小多项式的次数至多为 \(n\)，这是前面提到过的一个结果。

命题 \ref{pro:12.20} 中的关系在确定矩阵 \(A\) 的不变因子时常常相当有用，特别是对次数较小的矩阵（参见习题 3、4 与例子）。下面的结果（它依赖于上一节的习题 16 到 19，其证明在习题中概述）一般地计算不变因子。

设 \(A\) 是域 \(F\) 上的 \(n \times n\) 矩阵。则 \(xI-A\) 是元素在 \(F[x]\) 中的 \(n \times n\) 矩阵。三种操作

（a）交换两行或两列；

（b）把一行或一列的（\(F[x]\) 中的）倍数加到另一行或另一列上；

（c）把任意行或列乘以 \(F[x]\) 中的单位，即 \(F\) 中的非零元素，

称为\textbf{初等行运算与列运算}。

\begin{theorem}\label{thm:12.21}
设 \(A\) 是域 \(F\) 上的 \(n \times n\) 矩阵。用上面三种初等行运算与列运算，元素在 \(F[x]\) 中的 \(n \times n\) 矩阵 \(xI-A\) 可以化为对角形（称为 \(A\) 的\textbf{Smith 标准形}）
\[
\begin{pmatrix}
1 & & & & \\
& \ddots & & & \\
& & 1 & & \\
& & & a_1(x) & \\
& & & & \ddots & \\
& & & & & a_m(x)
\end{pmatrix}
\]
其中 \(a_1(x), a_2(x), \ldots, a_m(x)\) 是 \(F[x]\) 中首一非零元素，次数至少为 1，且满足 \(a_1(x) \mid a_2(x) \mid \cdots \mid a_m(x)\). 元素 \(a_1(x), \ldots, a_m(x)\) 就是 \(A\) 的不变因子。
\end{theorem}

\begin{proof}
参见习题。
\end{proof}

\noindent\textbf{不变因子分解算法：化矩阵为有理标准形}

如上一节末尾的习题中提到的，记录使 \(xI-A\) 对角化所需的运算将显式地给出一个矩阵 \(P\)，使得 \(P^{-1}AP\) 处于有理标准形。等价地，若 \(V\) 是给定的以向量空间基 \([e_1, e_2, \ldots, e_n]\) 的 \(F[x]\)-模，则 \(P\) 定义了给出 \(V\) 分解为循环 \(F[x]\)-模直和的不变因子分解的基变换。特别地，若 \(A\) 是由 \(x\) 定义的 \(F[x]\)-模 \(V\) 的线性变换 \(T\) 的矩阵（即 \(T(e_j) = xe_j = \sum_{i=1}^n a_{ij}e_i\)，其中 \(A = (a_{ij})\)），则矩阵 \(P\) 定义了使 \(T\) 的矩阵处于有理标准形的 \(V\) 的基变换。

我们首先在确定给定以向量空间基 \([e_1, e_2, \ldots, e_n]\) 的 \(F[x]\)-模 \(V\) 的不变因子分解的一般背景下描述这个算法（其证明在习题中概述）。然后我们描述把给定 \(n \times n\) 矩阵 \(A\) 化为有理标准形的算法（其中对底层向量空间与关联线性变换的指称被省略）。

\noindent\textbf{不变因子分解算法}

设 \(V\) 是以向量空间基 \([e_1, e_2, \ldots, e_n]\) 的 \(F[x]\)-模（故特别地这些元素作为 \(F[x]\)-模生成 \(V\)）。设 \(T\) 是由 \(x\) 定义的 \(V\) 到自身的线性变换，\(A\) 是与 \(T\) 及这组 \(V\) 的基关联的 \(n \times n\) 矩阵，即 \(T(e_j) = xe_j = \sum_{i=1}^n a_{ij}e_i\).

（1）用下列三种初等行运算与列运算把 \(xI-A\) 在 \(F[x]\) 上对角化，并记录所用的行运算：（a）交换两行或两列（交换第 \(i\) 行与第 \(j\) 行记作 \(R_i \leftrightarrow R_j\)，对列类似地记作 \(C_i \leftrightarrow C_j\)）；（b）把一行或一列的（\(F[x]\) 中的）倍数加到另一行或另一列上（若把第 \(j\) 行的 \(p(x)\) 倍加到第 \(i\) 行上记作 \(R_i+p(x)R_j \mapsto R_i\)，对列类似地记作 \(C_i+p(x)C_j \mapsto C_i\)）；（c）把任意行或列乘以 \(F[x]\) 中的单位，即 \(F\) 中的非零元素（若把第 \(i\) 行乘以 \(u \in F^\times\) 记作 \(uR_i\)，对列类似地记作 \(uC_i\)）。

（2）从 \(F[x]\)-模生成元 \([e_1, e_2, \ldots, e_n]\) 开始，对（1）中用到的每个行运算，按下列规则改变生成元组：（a）若交换第 \(i\) 行与第 \(j\) 行，则交换第 \(i\) 个与第 \(j\) 个生成元；（b）若把第 \(j\) 行的 \(p(x)\) 倍加到第 \(i\) 行上，则从第 \(j\) 个生成元中减去第 \(i\) 个生成元的 \(p(x)\) 倍（注意指标）；（c）若把第 \(i\) 行乘以单位 \(u \in F\)，则把第 \(i\) 个生成元除以 \(u\).

（3）当 \(xI-A\) 已对角化为定理 \ref{thm:12.21} 中的形式时，\(V\) 的生成元 \([e_1, e_2, \ldots, e_n]\) 将具有 \(e_1, e_2, \ldots, e_n\) 的 \(F[x]\)-线性组合的形式。用 \(xe_j = T(e_j) = \sum_{i=1}^n a_{ij}e_i\) 把这些元素写成 \(e_1, e_2, \ldots, e_n\) 的 \(F\)-线性组合。当 \(xI-A\) 已对角化时，这些线性组合中前 \(n-m\) 个为 0（这提供有用的数值检验），其余线性组合非零，即 \(V\) 的生成元具有 \([0, \ldots, 0, f_1, \ldots, f_m]\) 的形式，恰好与定理 \ref{thm:12.21} 中的对角元素对应。元素 \(f_1, \ldots, f_m\) 是 \(V\) 的不变因子分解中各个循环因子的一组 \(F[x]\)-模生成元（零化子分别为 \((a_1(x)), \ldots, (a_m(x))\)）：
\[
V = F[x]f_1 \oplus F[x]f_2 \oplus \cdots \oplus F[x]f_m, \qquad F[x]f_i \cong F[x]/(a_i(x)),\ i = 1, 2, \ldots, m,
\]
这给出 \(F[x]\)-模 \(V\) 的不变因子分解。

（4）\(V\) 的每个循环因子相应的向量空间基此时由 \(f_i, xf_i = Tf_i, x^2f_i = T^2f_i, \ldots, x^{d_i-1}f_i = T^{d_i-1}f_i\) 给出，其中 \(d_i\) 是 \(a_i(x)\) 的次数。

（5）把（4）中计算出的每个循环因子的向量空间基用原来的向量空间基 \([e_1, e_2, \ldots, e_n]\) 表示，并把这些坐标用作 \(n \times n\) 矩阵 \(P\) 的第 \(k\) 列。则 \(P^{-1}AP\) 处于有理标准形（对角块为 \(a_i(x)\) 的伴侣矩阵）。这就是 \(T\) 关于（4）中向量空间基的矩阵。

现在我们描述把给定 \(n \times n\) 矩阵 \(A\) 化为有理标准形（即确定使 \(P^{-1}AP\) 处于有理标准形的 \(n \times n\) 矩阵 \(P\)）的算法。这不过是把上述算法应用于向量空间 \(V = F^n\)（\(n \times 1\) 列向量以标准基 \([e_1, e_2, \ldots, e_n]\)，其中 \(e_i\) 是第 \(i\) 个位置为 1、其余位置为 0 的列向量），而 \(T\) 是由 \(A\) 及这组基定义的线性变换。省略对底层向量空间与关联线性变换的指称后，这个算法纯粹是矩阵论的：

\noindent\textbf{把 \(n \times n\) 矩阵化为有理标准形}

设 \(A\) 是元素在域 \(F\) 中的 \(n \times n\) 矩阵。

（1）用下列三种初等行运算与列运算把矩阵 \(xI-A\) 在 \(F[x]\) 上对角化，并记录所用的运算：（a）交换两行或两列（记作 \(R_i \leftrightarrow R_j\)、\(C_i \leftrightarrow C_j\)）；（b）把一行或一列的（\(F[x]\) 中的）倍数加到另一行或另一列上（记作 \(R_i+p(x)R_j \mapsto R_i\)、\(C_i+p(x)C_j \mapsto C_i\)）；（c）把任意行或列乘以 \(F[x]\) 中的单位，即 \(F\) 中的非零元素（记作 \(uR_i\)、\(uC_i\)）。定义 \(d_1, \ldots, d_m\) 为对角线上出现的首一非恒定多项式 \(a_1(x), \ldots, a_m(x)\) 的次数。

（2）从 \(n \times n\) 恒等矩阵 \(P'\) 开始，对（1）中用到的每个行运算，按下列规则改变矩阵 \(P'\)：（a）若 \(R_i \leftrightarrow R_j\)，则交换 \(P'\) 的第 \(i\) 列与第 \(j\) 列（即对 \(P'\) 作 \(C_i \leftrightarrow C_j\)）；（b）若 \(R_i+p(x)R_j \mapsto R_i\)，则从 \(P'\) 的第 \(j\) 列中减去矩阵 \(p(A)\) 乘以 \(P'\) 的**第 \(i\) 列**的乘积（即对 \(P'\) 作 \(C_j-p(A)C_i \mapsto C_j\)，注意指标）；（c）若 \(uR_i\)，则把 \(P'\) 的第 \(i\) 列的元素除以 \(u\)（即对 \(P'\) 作 \(u^{-1}C_i\)）。

（3）当 \(xI-A\) 已对角化为定理 \ref{thm:12.21} 中的形式时，矩阵 \(P'\) 的前 \(n-m\) 列为 0（这提供有用的数值检验），其余 \(m\) 列非零。对每个 \(i = 1, 2, \ldots, m\)，把 \(P'\) 的第 \(i\) 个非零列依次乘以 \(A^j\)（\(j = 0, 1, \ldots, d_i-1\)，其中 \(d_i\) 是上面（1）中的整数），并按此顺序把所得列向量用作 \(n \times n\) 矩阵 \(P\) 接下来的 \(d_i\) 列。则 \(P^{-1}AP\) 处于有理标准形（其对角块是（1）中多项式 \(a_1(x), \ldots, a_m(x)\) 的伴侣矩阵）。

在线性变换（或矩阵）标准形理论中，特征多项式扮演有限阿贝尔群的阶的角色，极小多项式扮演指数的角色（毕竟它们是同样的不变量，一个对应主理想整环 \(\mathbb{Z}\) 上的模，另一个对应主理想整环 \(F[x]\) 上的模），故我们可以直接解决与第 5 章中关于有限阿贝尔群考虑过的问题类似的问题。特别地，这包括：

（A）确定给定矩阵的有理标准形（类似于把有限阿贝尔群分解为循环群的直积）；

（B）判断两个给定矩阵是否相似（类似于判断两个给定有限阿贝尔群是否同构）；

（C）确定 \(F\) 上具有给定特征多项式的矩阵的全部相似类（类似于确定给定阶的全部阿贝尔群）；

（D）确定 \(F\) 上具有给定极小多项式的 \(n \times n\) 矩阵的全部相似类（类似于确定给定指数、秩至多为 \(n\) 的全部阿贝尔群）。

\begin{example}
（1）我们求下列 \(\mathbb{Q}\) 上矩阵的有理标准形并判断它们是否相似：
\[
A = \begin{pmatrix} 2 & -2 & 14 \\ 0 & 3 & -7 \\ 0 & 0 & 2 \end{pmatrix}, \quad B = \begin{pmatrix} 0 & -4 & 85 \\ 1 & 4 & -30 \\ 0 & 0 & 3 \end{pmatrix}, \quad C = \begin{pmatrix} 2 & 2 & 1 \\ 0 & 2 & -1 \\ 0 & 0 & 3 \end{pmatrix}.
\]
直接计算表明这三个矩阵有相同的特征多项式：\(c_A(x) = c_B(x) = c_C(x) = (x-2)^2(x-3)\). 由于极小多项式与特征多项式有相同的根，极小多项式的可能取值只有 \((x-2)(x-3)\) 或 \((x-2)^2(x-3)\). 快速计算得 \((A-2I)(A-3I) = 0\)、\((B-2I)(B-3I) \neq 0\)（其 \(1,1\) 元素非零）与 \((C-2I)(C-3I) \neq 0\)（其 \(1,2\) 元素非零）。由此可见
\[
m_A(x) = (x-2)(x-3), \qquad m_B(x) = m_C(x) = (x-2)^2(x-3).
\]
由此立即得知 \(B\) 与 \(C\) 没有额外的不变因子。由于 \(A\) 的不变因子整除极小多项式且其积为特征多项式，我们看到 \(A\) 的不变因子是多项式 \(x-2\) 与 \((x-2)(x-3) = x^2-5x+6\)（对 \(2 \times 2\) 与 \(3 \times 3\) 矩阵，特征多项式与极小多项式的确定给出全部不变因子，参见习题 3、4）。我们断定 \(B\) 与 \(C\) 相似，而二者都不与 \(A\) 相似。它们的有理标准形是（注意 \((x-2)^2(x-3) = x^3-7x^2+16x-12\)）
\[
\begin{pmatrix} 2 & 0 & 0 \\ 0 & 0 & 6 \\ 0 & 1 & 5 \end{pmatrix}, \quad \begin{pmatrix} 0 & 0 & 12 \\ 1 & 0 & -16 \\ 0 & 1 & 7 \end{pmatrix}, \quad \begin{pmatrix} 0 & 0 & 12 \\ 1 & 0 & -16 \\ 0 & 1 & 7 \end{pmatrix}.
\]

（2）上一例中的有理标准形只是通过确定矩阵的特征多项式与极小多项式得到的。如前所述，对 \(2 \times 2\) 与 \(3 \times 3\) 矩阵这已经足够，因为这些信息足以确定全部不变因子。然而对更大的矩阵，这通常不够（参见下一例），需要更多工作来确定不变因子。本例如上一例中的矩阵 \(A\) 那样，按上面概述的两个算法再次计算有理标准形。虽然对这个很小的矩阵这在计算上更困难（这一点将会显现），但它的优点是即使在这种情况下也能显式地计算出使 \(P^{-1}AP\) 处于有理标准形的矩阵 \(P\).

\noindent I.（不变因子分解）我们用 \(\mathbb{Q}[x]\) 中的行运算与列运算把矩阵
\[
xI-A = \begin{pmatrix} x-2 & 2 & -14 \\ 0 & x-3 & 7 \\ 0 & 0 & x-2 \end{pmatrix}
\]
化为对角形。与不变因子分解算法中一样，我们用记号 \(R_i \leftrightarrow R_j\) 表示交换第 \(i\) 行与第 \(j\) 行，\(R_i+aR_j \mapsto R_i\) 表示把第 \(j\) 行的 \(a\) 倍加到第 \(i\) 行上，\(uR_i\) 表示把第 \(i\) 行乘以 \(u\)（列的情形类似，用 \(C\) 代替 \(R\)）。注意下面施行的前两个运算相当特别，只是为了在计算中处处得到整数：
\[
\begin{pmatrix} x-2 & 2 & -14 \\ 0 & x-3 & 7 \\ 0 & 0 & x-2 \end{pmatrix} \longrightarrow \begin{pmatrix} x-2 & x-1 & -7 \\ 0 & x-3 & 7 \\ 0 & 0 & x-2 \end{pmatrix} \longrightarrow \begin{pmatrix} -1 & x-1 & -x+1 \\ -x+3 & x-3 & x-3 \\ 0 & 0 & 0 \end{pmatrix}
\]
（此处按原文依次施用 \(C_2+C_1 \mapsto C_2\) 等列运算与行运算；完整的运算序列为
\[
C_2+C_1 \mapsto C_2,\quad R_1 \leftrightarrow R_2,\quad -R_1,\quad R_2+(x-3)R_1 \mapsto R_2,\quad C_2+(x-1)C_1 \mapsto C_2,
\]
\[
C_3 \mapsto C_3-C_2,\quad R_2+(x-2)R_1 \mapsto R_2,\quad R_3 \leftrightarrow R_2,\quad C_2 \leftrightarrow C_3,
\]
最后得到对角矩阵 \(\operatorname{diag}(x-2,\ x^2-5x+6,\ 0)\)）。

这确定了这个矩阵的不变因子 \(x-2\) 与 \(x^2-5x+6\)，即我们在上面例 1 中确定的不变因子。现在设 \(V\) 是 \(\mathbb{Q}\) 上以 \(e_1, e_2, e_3\) 为基的 3 维向量空间，\(T\) 是相应的线性变换（它定义 \(x\) 在 \(V\) 上的作用），即 \(xe_1 = T(e_1) = 2e_1\)、\(xe_2 = T(e_2) = -2e_1+3e_2\)、\(xe_3 = T(e_3) = 14e_1-7e_2+2e_3\). 上面约化中用到的行运算为 \(R_1 \leftrightarrow R_2\)、\(-R_1\)、\(R_2+(x-3)R_1 \mapsto R_2\)、\(R_2-7R_3 \mapsto R_2\)、\(R_2 \leftrightarrow R_3\).

从 \(V\) 的基 \([e_1, e_2, e_3]\) 出发，按正文中给出的规则改变它，我们得到
\[
[e_1, e_2, e_3] \to [e_1, e_2-e_1, e_3] \to [-e_1, e_2-e_1, e_3] \to [-e_1-(x-3)(e_2-e_1), e_2-e_1, e_3]
\]
\[
\to [-e_1-(x-3)(e_2-e_1), e_2-e_1, e_3+7(e_2-e_1)] \to [-e_1-(x-3)(e_2-e_1), e_3+7(e_2-e_1), e_2-e_1].
\]
利用上面 \(x\) 的作用公式，我们看到最后这些元素是 \(V\) 中对应于 \(xI-A\) 对角化形式中元素 \(1\)、\(x-2\) 与 \(x^2-5x+6\) 的元素 \([0, -7e_1+7e_2+e_3, -e_1+e_2]\). 因此元素 \(f_1 = -7e_1+7e_2+e_3\) 与 \(f_2 = -e_1+e_2\) 是 \(V\) 作为 \(\mathbb{Q}[x]\)-模的不变因子分解中两个循环因子的 \(\mathbb{Q}[x]\)-模生成元。这两个因子相应的 \(\mathbb{Q}\)-向量空间基是 \(f_1\) 与 \(f_2,\ xf_2 = Tf_2\)，即 \(-7e_1+7e_2+e_3\) 与 \(-e_1+e_2\)、\(T(-e_1+e_2) = -4e_1+3e_2\). 于是矩阵
\[
P = \begin{pmatrix} -7 & -1 & -4 \\ 7 & 1 & 3 \\ 1 & 0 & 0 \end{pmatrix}
\]
把 \(A\) 共轭为它的有理标准形
\[
P^{-1}AP = \begin{pmatrix} 2 & 0 & 0 \\ 0 & 0 & -6 \\ 0 & 1 & 5 \end{pmatrix},
\]
容易验证。

\noindent II.（直接把 \(A\) 化为有理标准形）我们用 \(xI-A\) 对角化中用到的行运算按上面的算法确定矩阵 \(P'\)：
\[
P' = \begin{pmatrix} -7 & -1 & 0 \\ 7 & 1 & 0 \\ 0 & 0 & 1 \end{pmatrix}.
\]
这里 \(d_1 = 1\)、\(d_2 = 2\)，分别对应于 \(P'\) 的第二个与第三个非零列。故 \(P\) 的各列分别由
\[
\begin{pmatrix} -1 \\ 1 \\ 0 \end{pmatrix},\quad \begin{pmatrix} -7 \\ 7 \\ 1 \end{pmatrix},\quad A\begin{pmatrix} -1 \\ 1 \\ 0 \end{pmatrix} = \begin{pmatrix} -4 \\ 3 \\ 0 \end{pmatrix}
\]
给出，再次得到上面的矩阵 \(P\).

（3）对 \(3 \times 3\) 矩阵 \(A\)，如上面例 1 所见，不必施行上述冗长计算就确定了有理标准形（等价地，不变因子）。然而对 \(n \times n\) 且 \(n \geq 4\) 的矩阵，特征多项式与极小多项式的计算一般不足以确定全部不变因子，故上一例中更繁复的计算可能是必要的。例如考虑矩阵
\[
D = \begin{pmatrix} 1 & -1 & 2 & -1 \\ 1 & 0 & 1 & -2 \\ 0 & 1 & -2 & 3 \end{pmatrix}
\]
（此处按原文，\(D\) 是 \(4 \times 4\) 矩阵 \(\begin{pmatrix} 1 & -1 & 2 & -1 \\ 1 & 0 & 1 & -2 \\ 1 & 0 & 1 & -2 \\ 0 & 1 & -2 & 3\end{pmatrix}\)）。

简短计算表明 \(D\) 的特征多项式为 \((x-1)^4\). 可能的极小多项式为 \(x-1\)、\((x-1)^2\)、\((x-1)^3\) 与 \((x-1)^4\). 显然 \(D-I \neq 0\)，而另一次简短计算表明 \((D-I)^2 = 0\)，故 \(D\) 的极小多项式为 \((x-1)^2\). 于是不变因子有两种可能：\(x-1, x-1, (x-1)^2\) 与 \((x-1)^2, (x-1)^2\).

为确定 \(D\) 的不变因子，我们把上一例的步骤应用于 \(4 \times 4\) 矩阵
\[
xI-D = \begin{pmatrix} x-1 & 1 & -2 & 1 \\ -1 & x & -1 & 2 \\ -1 & 0 & x-1 & 2 \\ 0 & -1 & 2 & x-3 \end{pmatrix}.
\]
由初等行运算与列运算从它得到的对角矩阵是
\[
\begin{pmatrix} 1 & 0 & 0 & 0 \\ 0 & 1 & 0 & 0 \\ 0 & 0 & (x-1)^2 & 0 \\ 0 & 0 & 0 & (x-1)^2 \end{pmatrix},
\]
这表明 \(D\) 的不变因子是 \((x-1)^2, (x-1)^2\).

\noindent I.（不变因子分解）若此时 \(e_1, e_2, e_3, e_4\) 是 \(V\) 的一组基，则如上一例那样使用对角化中的行运算，我们看到对应于上述因子的 \(V\) 的生成元是 \((x-1)e_1-2e_2-e_3 = 0\)、\(-2e_1+(x+1)e_2-e_4 = 0\)、\(e_1\)、\(e_2\). 因此此时 \(V\) 不变因子分解中两个直因子的一组向量空间基由 \(e_1, Te_1\) 与 \(e_2, Te_2\) 给出，其中 \(T\) 是由 \(D\) 定义的线性变换，即 \(e_1\)、\(e_1+2e_2+e_3\) 与 \(e_2\)、\(2e_1-e_2+e_4\). 联系这些基的相应矩阵 \(P\) 是
\[
P = \begin{pmatrix} 1 & 1 & 0 & 2 \\ 0 & 2 & 1 & -1 \\ 0 & 1 & 0 & 0 \\ 0 & 0 & 0 & 1 \end{pmatrix},
\]
使得 \(P^{-1}DP\) 处于有理标准形
\[
P^{-1}DP = \begin{pmatrix} 0 & -1 & 0 & 0 \\ 1 & 2 & 0 & 0 \\ 0 & 0 & 0 & -1 \\ 0 & 0 & 1 & 2 \end{pmatrix},
\]
容易验证。

\noindent II.（直接把 \(D\) 化为有理标准形）如例 2 那样，我们从 \(xI-D\) 对角化中用到的行运算确定算法的矩阵 \(P'\)，得到
\[
P' = \begin{pmatrix} 1 & 0 & 0 & 0 \\ 0 & 1 & 0 & 0 \\ 0 & 0 & 0 & 1 \\ 0 & 0 & 1 & 0 \end{pmatrix},
\]
其中 \(d_1 = 2\)、\(d_2 = 2\)，分别对应于 \(P'\) 的第三与第四个非零列。故 \(P\) 的各列由此再次得到上面的矩阵 \(P\).

（4）本例我们确定元素取自 \(\mathbb{Q}\)、特征多项式为 \((x^4-1)(x^2-1)\) 的矩阵 \(A\) 的全部相似类。首先注意特征多项式次数为 6 的矩阵必是 \(6 \times 6\) 矩阵。多项式 \((x^4-1)(x^2-1)\) 在 \(\mathbb{Q}[x]\) 中分解为不可约因子之积：\((x-1)^2(x+1)^2(x^2+1)\). 由于 \(A\) 的极小多项式 \(m_A(x)\) 与 \(c_A(x)\) 有相同的根，可知 \((x-1)(x+1)(x^2+1)\) 整除 \(m_A(x)\). 假设 \(a_1(x), \ldots, a_m(x)\) 是某个 \(A\) 的不变因子，则 \(a_m(x) = m_A(x)\)，\(a_i(x) \mid a_{i+1}(x)\)（特别地所有不变因子都整除 \(m_A(x)\)），且 \(a_1(x)a_2(x)\cdots a_m(x) = (x^4-1)(x^2-1)\). 容易看出在这些约束下唯一允许的表是

（a）\((x-1)(x+1),\ (x-1)(x+1)(x^2+1)\)；

（b）\(x-1,\ (x-1)(x+1)^2(x^2+1)\)；

（c）\(x+1,\ (x-1)^2(x+1)(x^2+1)\)；

（d）\((x-1)^2(x+1)^2(x^2+1)\).

现在容易写出相应的伴侣矩阵的直和，得到这 4 个相似类的代表元。我们将在下一节看到，即使在 \(M_6(\mathbb{C})\) 中仍然只有这 4 个相似类。

（5）本例我们求元素取自 \(\mathbb{Q}\)、满足 \(A^6 = I\) 的 \(3 \times 3\) 矩阵 \(A\) 的相似类。对每个这样的 \(A\)，其极小多项式整除 \(x^6-1\)，而这个多项式在 \(\mathbb{Q}[x]\) 中的完全分解为
\[
x^6-1 = (x-1)(x+1)(x^2-x+1)(x^2+x+1).
\]
反之，若 \(B\) 是极小多项式整除 \(x^6-1\) 的任意 \(3 \times 3\) 矩阵，则 \(B^6 = I\). 对 \(B\) 的极小多项式的唯一限制是其次数至多为 3（由凯莱--哈密顿定理）。故这样的矩阵 \(A\) 的极小多项式的可能取值是

（a）\(x-1\)；（b）\(x+1\)；（c）\(x^2-x+1\)；（d）\(x^2+x+1\)；（e）\((x-1)(x+1)\)；（f）\((x-1)(x^2-x+1)\)；（g）\((x-1)(x^2+x+1)\)；（h）\((x+1)(x^2-x+1)\)；（i）\((x+1)(x^2+x+1)\).

在有理标准形的约束下，这些给出下列允许的不变因子表：

（i）\(x-1, x-1, x-1\)；（ii）\(x+1, x+1, x+1\)；（iii）\(x-1, (x-1)(x+1)\)；（iv）\(x+1, (x-1)(x+1)\)；（v）\((x-1)(x^2-x+1)\)；（vi）\((x-1)(x^2+x+1)\)；（vii）\((x+1)(x^2-x+1)\)；（viii）\((x+1)(x^2+x+1)\).

注意若极小多项式是 \(x^2+x+1\) 或 \(x^2-x+1\)，则不可能有合适的不变因子组。现在可以写出相应的有理标准形；例如（i）是 \(I\)，（ii）是 \(-I\)，而（iii）是
\[
\begin{pmatrix} 1 & 0 & 0 \\ 0 & 0 & 1 \\ 0 & 1 & 0 \end{pmatrix}.
\]
还要注意这个结果的另一种表述是：元素取自 \(\mathbb{Q}\)、其（当然是乘法的）阶整除 6 的任何 \(3 \times 3\) 矩阵都相似于这 8 个矩阵之一，故本例确定了 \(GL_3(\mathbb{Q})\) 中阶为 1、2、3 与 6 的全部元素（在相似意义下）。
\end{example}

\subsection*{习题}

\begin{enumerate}
\item 证明 \(V\) 的相似线性变换（或相似的 \(n \times n\) 矩阵）有相同的特征多项式与相同的极小多项式。

\item 设 \(M\) 如引理 \ref{lem:12.19} 中那样。证明 \(M\) 的极小多项式是 \(A_1, \ldots, A_k\) 的极小多项式的最小公倍。

\item 证明 \(F\) 上两个不是纯量矩阵的 \(2 \times 2\) 矩阵相似当且仅当它们有相同的特征多项式。

\item 证明两个 \(3 \times 3\) 矩阵相似当且仅当它们有相同的特征多项式与相同的极小多项式。对 \(4 \times 4\) 矩阵给出这一断言的一个明确反例。

\item 由 \(F\) 上 \(n\) 维向量空间 \(V\) 的所有线性变换本身构成 \(F\) 上维数为 \(n^2\) 的向量空间这一事实，直接证明线性变换 \(T\) 的极小多项式的次数至多为 \(n^2\).

\item 证明 \(n \times n\) 矩阵 \(A\) 的特征多项式的常数项是 \((-1)^n\det A\)，而 \(x^{n-1}\) 的系数是对角元素之和的相反数（\(A\) 的对角元素之和称为 \(A\) 的\textbf{迹}）。证明 \(\det A\) 是 \(A\) 的特征值之积，而 \(A\) 的迹是 \(A\) 的特征值之和。

\item 确定矩阵
\[
\begin{pmatrix} 0 & 0 & 0 \\ 1 & 0 & 0 \\ 0 & 1 & 0 \end{pmatrix}
\]
的特征值。（按原文，此矩阵为 \(4 \times 4\) 的 \(\begin{pmatrix}0&0&0&0\\1&0&0&0\\0&1&0&0\\0&0&1&0\end{pmatrix}\)——伴侣矩阵 \(C_{x^4}\).）

\item 验证伴侣矩阵
\[
\begin{pmatrix} 0 & 0 & \cdots & 0 & -a_0 \\ 1 & 0 & \cdots & 0 & -a_1 \\ 0 & 1 & \cdots & 0 & -a_2 \\ \vdots & & \ddots & & \vdots \\ 0 & 0 & \cdots & 1 & -a_{k-1} \end{pmatrix}
\]
的特征多项式是 \(x^k+a_{k-1}x^{k-1}+\cdots+a_1x+a_0\).

\item 求
\[
\begin{pmatrix} 0 & -1 & -1 \\ 0 & 0 & 0 \\ -1 & 0 & 0 \end{pmatrix}, \quad \begin{pmatrix} c & 0 & -1 \\ 0 & c & 1 \\ 1 & 0 & c \end{pmatrix}, \quad \begin{pmatrix} 422 & 465 & 15 & -30 \\ -420 & -463 & -15 & 30 \\ 840 & 930 & 32 & -60 \\ -140 & -155 & -5 & 12 \end{pmatrix}
\]
的有理标准形。

\item 求 \(\mathbb{Q}\) 上极小多项式为 \((x+2)^2(x-1)\) 的 \(6 \times 6\) 矩阵的全部相似类（只需给出全部不变因子表并写出其中一些相应的矩阵）。

\item 求 \(\mathbb{C}\) 上特征多项式为 \((x^4-1)(x^2-1)\) 的 \(6 \times 6\) 矩阵的全部相似类。

\item 求 \(\mathbb{F}_2\) 上满足 \(A^6 = I\) 的 \(3 \times 3\) 矩阵 \(A\) 的全部相似类（与我们在 \(\mathbb{Q}\) 上算出的答案作比较）。对满足 \(B^{20} = I\) 的 \(4 \times 4\) 矩阵 \(B\) 做同样的事。

\item 证明 \(\mathbb{Q}\) 上具有给定特征多项式（在 \(\mathbb{Q}[x]\) 中）的 \(3 \times 3\) 矩阵的相似类个数与在 \(\mathbb{Q}\) 的任何扩域上的相似类个数相同。给出一个例子表明对 \(4 \times 4\) 矩阵这一断言一般不再成立。

\item 确定特征多项式为 \(x^2(x^2+1)^2\) 的线性变换的全部可能的有理标准形。

\item 在相似意义下确定所有（乘法）阶恰为 4 的 \(2 \times 2\) 有理矩阵（即 \(\in M_2(\mathbb{Q})\)）。若矩阵的元素取自 \(\mathbb{C}\) 也做同样的事。

\item 证明在 \(\mathbb{F}_{19}[x]\) 中 \(x^5-1 = (x-1)(x^2-4x+1)(x^2+5x+1)\). 用它确定元素取自 \(\mathbb{F}_{19}\)、（乘法）阶为 5 的全部 \(2 \times 2\) 矩阵（在相似意义下）。

\item 确定 \(GL_3(\mathbb{F}_2)\) 的共轭类的代表元。〔把你的答案与第 6 章定理 15 和命题 14 作比较。〕

\item 设 \(V\) 是 \(\mathbb{Q}\) 上的有限维向量空间，并假设 \(T\) 是 \(V\) 的非奇异线性变换使得 \(T^{-1} = T^2+T\). 证明 \(V\) 的维数被 3 整除。若 \(V\) 的维数恰为 3，证明所有这样的变换 \(T\) 都相似。

\item 设 \(V\) 是无限维实向量空间 \(\mathbb{R}^\infty = \{(a_0, a_1, a_2, \ldots) \mid a_0, a_1, a_2, \ldots \in \mathbb{R}\}\). 由 \(T(a_0, a_1, a_2, \ldots) = (0, a_0, a_1, a_2, \ldots)\) 定义映射 \(T : V \to V\). 证明 \(T\) 没有特征向量。

\item 设 \(\ell\) 是素数，\(\Phi_\ell(x) = \frac{x^\ell-1}{x-1} = x^{\ell-1}+x^{\ell-2}+\cdots+x+1 \in \mathbb{Z}[x]\) 是第 \(\ell\) 个分圆多项式，它在 \(\mathbb{Q}\) 上不可约（第 9.4 节推论 \ref{cor:9.14} 之后的例 4）。本习题确定 \(\Phi_\ell(x)\) 模 \(p\) 的因子的最小次数，从而特别地确定 \(\Phi_\ell(x)\) 模 \(p\) 何时不可约。（这实际上确定了 \(\Phi_\ell(x)\) 模 \(p\) 的完全分解——参见第 13.6 节习题 8。）

（a）证明若 \(p = \ell\)，则 \(\Phi_\ell(x)\) 在 \(\mathbb{F}_\ell[x]\) 中被 \(x-1\) 整除。

（b）假设 \(p \neq \ell\)，用 \(f\) 表示 \(p\) 在 \(\mathbb{F}_\ell^\times\) 中的阶，即 \(f\) 是使 \(p^f \equiv 1 \bmod \ell\) 的最小 \(p\) 的幂。证明 \(m = f\) 是使群 \(GL_m(\mathbb{F}_p)\) 含有阶为 \(\ell\) 的元素 \(A\) 的最小 \(m\) 值。〔利用第 11.1 节末尾 \(GL(V)\) 的阶的公式。〕

（c）证明 \(\Phi_\ell(x)\) 在 \(\mathbb{F}_p[x]\) 中不被任何次数小于 \(f\) 的多项式整除〔考虑这样因子的伴侣矩阵并利用（b）〕。设 \(m_A(x) \in \mathbb{F}_p[x]\) 表示（b）中矩阵 \(A\) 的极小多项式，并断定 \(m_A(x)\) 是不可约的、次数为 \(f\)、且在 \(\mathbb{F}_p[x]\) 中整除 \(\Phi_\ell(x)\).

（d）特别地证明 \(\Phi_\ell(x)\) 模 \(p\) 不可约当且仅当 \(f = \ell-1\) 是使 \(p\) 的幂同余于 1 模 \(\ell\) 的最小幂，即 \(p\) 是模 \(\ell\) 的原根。

\item 证明定理 \ref{thm:12.21} 之前描述的三种初等行运算与列运算中前两种不改变矩阵的行列式，第三种把行列式乘以一个单位。由定理 \ref{thm:12.21} 断定 \(A\) 的特征多项式与 \(A\) 的不变因子之积相差一个单位。由于按定义这两个多项式都是首一的，断定它们相等（这给出命题 \ref{pro:12.20} 的另一个证明）。
\end{enumerate}

\noindent 下面各习题概述定理 \ref{thm:12.21} 的证明。它们对欧几里得整环 \(F[x]\) 显式地实现了上一节习题 16 到 19 中描述的构造。设 \(V\) 是以 \(v_1, v_2, \ldots, v_n\) 为基的 \(n\) 维向量空间，\(T\) 是由矩阵 \(A\) 及这组基定义的 \(V\) 的线性变换，即 \(T(v_j) = \sum_{i=1}^n a_{ij}v_i\)（\(j = 1, 2, \ldots, n\)），其中 \(A = (a_{ij})\). 设 \(F[x]^n\) 是 \(F[x]\) 上秩 \(n\) 的自由模，\(\xi_1, \xi_2, \ldots, \xi_n\) 是一组基。则我们有把 \(\xi_i\) 映到 \(v_i\)（\(i = 1, 2, \ldots, n\)）的自然满 \(F[x]\)-模同态 \(\varphi : F[x]^n \to V\). 如上一节习题中所指出的，一旦确定了 \(\ker\varphi\) 的一组生成元与相应关系矩阵，就可以确定 \(F[x]\)-模 \(V\) 的不变因子。由于按定义 \(x\) 通过线性变换 \(T\) 作用于 \(V\)，我们有 \(x(v_j) = \sum_{i=1}^n a_{ij}v_i\)（\(j = 1, 2, \ldots, n\)）。

\begin{enumerate}
\setcounter{enumi}{21}
\item 证明元素
\[
\nu_j = -a_{1j}\xi_1-\cdots-a_{j-1,j}\xi_{j-1}+(x-a_{jj})\xi_j-a_{j+1,j}\xi_{j+1}-\cdots-a_{nj}\xi_n
\]
（\(j = 1, 2, \ldots, n\)）是 \(\ker\varphi\) 的元素。

\item （a）证明 \(x\xi_1 = \nu_1+\vartheta\)，其中 \(\vartheta \in F\xi_1+\cdots+F\xi_n\) 是由 \(\xi_1, \ldots, \xi_n\) 张成的 \(F\)-向量空间中的元素。

（b）证明 \(F[x]\xi_1+\cdots+F[x]\xi_n = (F[x]\nu_1+\cdots+F[x]\nu_n)+(F\xi_1+\cdots+F\xi_n)\).

\item 证明 \(\nu_1, \ldots, \nu_n\) 生成 \(\ker\varphi\).〔用上一个结果证明 \(\ker\varphi\) 的任何元素都是 \(\nu_1, \nu_2, \ldots, \nu_n\) 生成的模中元素与形如 \(b_1\xi_1+\cdots+b_n\xi_n\)（其中 \(b_i \in F\)）的元素之和。然后证明这样的元素在 \(\ker\varphi\) 中当且仅当所有 \(b_i\) 都为 0，因为 \(v_1, \ldots, v_n\) 是 \(V\) 在 \(F\) 上的基。〕

\item 证明 \(\ker\varphi\) 的生成元 \(\nu_1, \nu_2, \ldots, \nu_n\) 的相应关系矩阵是
\[
\begin{pmatrix} x-a_{11} & -a_{12} & \cdots & -a_{1n} \\ -a_{21} & x-a_{22} & \cdots & -a_{2n} \\ \vdots & & \ddots & \vdots \\ -a_{n1} & -a_{n2} & \cdots & x-a_{nn} \end{pmatrix} = xI-A^t,
\]
其中 \(A^t\) 是 \(A\) 的转置。断定定理 \ref{thm:12.21} 与确定 \(A\) 的不变因子的算法由上一节习题 16 到 19 推出（注意为使这个关系矩阵对角化所需的行运算与列运算正是使定理 \ref{thm:12.21} 中矩阵对角化所需的列运算与行运算，这解释了为什么不变因子算法要记录所用的行运算）。
\end{enumerate}

\section{若尔当标准形}

我们沿用上一节的记号：\(F\) 是域，\(F[x]\) 是系数在 \(F\) 中的 \(x\) 的多项式环，\(V\) 是 \(F\) 上 \(n\) 维有限维向量空间，\(T\) 是 \(V\) 的固定线性变换，我们经由它把 \(V\) 变成 \(F[x]\)-模，\(A\) 是系数在 \(F\) 中的 \(n \times n\) 矩阵。回忆一旦固定了 \(V\) 的一组基，任何线性变换 \(T\) 就定义一个矩阵 \(A\)，反之任何矩阵 \(A\) 也定义一个线性变换 \(T\).

上一节我们用主理想整环 \(F[x]\) 上有限生成模基本定理的不变因子形式得到了线性变换 \(T\) 的有理标准形与 \(n \times n\) 矩阵 \(A\) 的有理标准形。本节我们用基本定理的初等因子形式得到若尔当标准形。我们将看到，处于这个标准形的矩阵尽可能接近对角矩阵，故这些矩阵比有理标准形中的更简单（但我们失去了一些“有理性”的结果）。

模的初等因子是其不变因子的素幂因子（这是推论 \ref{cor:12.10}）。对 \(F[x]\)-模 \(V\)，不变因子是次数至少为 1 的首一多项式 \(a_1(x), a_2(x), \ldots, a_m(x)\)（且 \(a_1(x) \mid a_2(x) \mid \cdots \mid a_m(x)\)），故相应的初等因子是这些多项式的不可约因子的幂。这些多项式只在乘以单位的意义下确定，而如不变因子的情形一样，我们可以要求它们首一来唯一地确定它们。

为得到最简单的初等因子，我们将假设多项式 \(a_1(x), a_2(x), \ldots, a_m(x)\) 完全分解为线性因子，即 \(V\) 的初等因子是一次多项式 \((x-\lambda)^k\) 的幂。由于初等因子之积是特征多项式，这等价于假设域 \(F\) 包含线性变换 \(T\)（等价地，表示 \(T\) 的矩阵 \(A\)）的所有特征值。

在关于 \(F\) 的这个假设下，由定理 \ref{thm:12.6} 立即推出 \(V\) 是形如 \(F[x]/(x-\lambda)^k\)（\(\lambda \in F\) 是 \(T\) 的一个特征值）的有限多个循环 \(F[x]\)-模的直和，这与 \(V\) 的初等因子相对应。

现在我们为对应于 \(V\) 的初等因子的每个直和项选取一组向量空间基，使得相应的 \(T\) 的矩阵特别简单。回忆由 \(F[x]\)-模结构的定义，作用在 \(V\) 上的线性变换 \(T\) 就是元素 \(x\) 在每个直和项 \(F[x]/(x-\lambda)^k\) 上由乘法作用。

考虑商 \(F[x]/(x-\lambda)^k\) 中的元素
\[
(x-\lambda)^{k-1}, \quad (x-\lambda)^{k-2}, \quad \ldots, \quad x-\lambda, \quad 1.
\]
把这些多项式中的每一个按 \(x\) 展开，我们看到把这些元素与 \(F[x]/(x-\lambda)^k\) 的 \(F\)-基 \(\overline{x}^{k-1}, \overline{x}^{k-2}, \ldots, \overline{x}, \overline{1}\) 联系起来的矩阵是上三角的、对角线上全为 1. 由于这个矩阵可逆（行列式为 1），可知上面的元素是 \(F[x]/(x-\lambda)^k\) 的一组 \(F\)-基。

关于这组基，乘以 \(x\) 的线性变换作用方式特别简单（注意 \(x = \lambda+(x-\lambda)\)，且在商中 \(\overline{(x-\lambda)^k} = 0\)）：
\[
\overline{(x-\lambda)^{k-1}} \mapsto \overline{(x-\lambda)^{k-2}}\cdot\overline{x} = \overline{(x-\lambda)^{k-2}}(\lambda+(x-\lambda)) = \lambda\overline{(x-\lambda)^{k-2}}+\overline{(x-\lambda)^{k-1}},
\]
\[
\overline{(x-\lambda)^{k-2}} \mapsto \lambda\overline{(x-\lambda)^{k-3}}+\overline{(x-\lambda)^{k-2}}, \qquad \ldots, \qquad \overline{x-\lambda} \mapsto \lambda \cdot \overline{1}+\overline{(x-\lambda)}, \qquad \overline{1} \mapsto \lambda\overline{1}.
\]
关于这组基，乘以 \(x\) 的矩阵因此是
\[
\begin{pmatrix}
\lambda & 1 & & & \\
& \lambda & 1 & & \\
& & \ddots & \ddots & \\
& & & \lambda & 1\\
& & & & \lambda
\end{pmatrix},
\]
其中空白处全为零。这样的矩阵有一个名称：

\begin{definition}
主对角线上全为 \(\lambda\)、第一条“上对角线”上全为 1（如上所示）的 \(k \times k\) 矩阵称为特征值为 \(\lambda\) 的 \textbf{\(k \times k\) 初等若尔当矩阵}，或特征值为 \(\lambda\) 的\textbf{大小为 \(k\) 的若尔当块}。
\end{definition}

把这一点应用于 \(V\) 在初等因子分解中的每个循环因子，我们得到 \(V\) 的一组向量空间基，关于它线性变换 \(T\) 的矩阵是 \(V\) 的初等因子相应的若尔当块的直和，即对角线上为若尔当块的分块对角矩阵。

注意这个矩阵在沿对角线置换各块的顺序的意义下由 \(F[x]\)-模 \(V\) 的初等因子唯一确定；反之由定理 \ref{thm:12.9}，初等因子表在同构意义下唯一确定模 \(V\)（作为 \(F[x]\)-模）。

\begin{definition}
（1）若一个矩阵是对角线上为若尔当块的分块对角矩阵，则称它处于\textbf{若尔当标准形}。

（2）线性变换 \(T\) 的\textbf{若尔当标准形}是表示 \(T\) 的、处于若尔当标准形的矩阵。
\end{definition}

我们已经证明任何线性变换 \(T\) 都有若尔当标准形。与有理标准形的情形一样，由初等因子的唯一性可知若尔当标准形在沿对角线置换若尔当块的意义下唯一（因此称为 \(T\) 的若尔当标准形）。我们把它总结为下面的定理。

\begin{theorem}\label{thm:12.22}
（\textbf{线性变换的若尔当标准形}）设 \(V\) 是域 \(F\) 上的有限维向量空间，\(T\) 是 \(V\) 的线性变换。假设 \(F\) 包含 \(T\) 的所有特征值。

（1）存在 \(V\) 的一组基，使得 \(T\) 关于它的矩阵处于若尔当标准形，即以 \(V\) 的初等因子相应的若尔当块为对角块的分块对角矩阵。

（2）\(T\) 的若尔当标准形在沿对角线置换若尔当块的意义下唯一。
\end{theorem}

与有理标准形一样，下面的定理给出 \(F\) 上 \(n \times n\) 矩阵的相应陈述。

\begin{theorem}\label{thm:12.23}
（\textbf{矩阵的若尔当标准形}）设 \(A\) 是域 \(F\) 上的 \(n \times n\) 矩阵，并假设 \(F\) 包含 \(A\) 的所有特征值。

（1）矩阵 \(A\) 相似于一个处于若尔当标准形的矩阵，即存在 \(F\) 上可逆的 \(n \times n\) 矩阵 \(P\) 使 \(P^{-1}AP\) 是以 \(A\) 的初等因子相应的若尔当块为对角块的分块对角矩阵。

（2）\(A\) 的若尔当标准形在沿对角线置换若尔当块的意义下唯一。
\end{theorem}

若尔当标准形与对角矩阵的差别只在于第一条上对角线上可能出现的 1（而且只有当有大小大于 1 的若尔当块时才出现），故它接近对角矩阵。下面这个结果特别表明矩阵 \(A\) 的若尔当标准形是尽可能接近对角矩阵的。

\begin{corollary}\label{cor:12.24}
（1）若矩阵 \(A\) 相似于对角矩阵 \(D\)，则 \(D\) 就是 \(A\) 的若尔当标准形。

（2）两个对角矩阵相似当且仅当它们的对角元素在置换意义下相同。
\end{corollary}

\begin{proof}
第一个断言由若尔当标准形的唯一性立即推出，因为对角矩阵本身处于若尔当形（若尔当块大小都为 1）。若尔当标准形的唯一性给出（2）。
\end{proof}

下一个推论给出判断矩阵 \(A\) 何时可以对角化的判别法。

\begin{corollary}\label{cor:12.25}
若 \(A\) 是元素取自 \(F\) 的 \(n \times n\) 矩阵且 \(F\) 包含 \(A\) 的所有特征值，则 \(A\) 在 \(F\) 上相似于对角矩阵当且仅当 \(A\) 的极小多项式没有重根。
\end{corollary}

\begin{proof}
假设 \(A\) 相似于对角矩阵。对角矩阵的极小多项式没有重根（它的根正是对角线上互异的元素）。由于相似矩阵有相同的极小多项式，可知 \(A\) 的极小多项式没有重根。

反之，假设 \(A\) 的极小多项式没有重根，设 \(B\) 是 \(A\) 的若尔当标准形。矩阵 \(B\) 是对角线上为初等若尔当矩阵的分块对角矩阵。由上一节末尾的习题，\(B\) 的极小多项式是各若尔当块的极小多项式的最小公倍。容易直接看出特征值为 \(\lambda\) 的大小为 \(k\) 的若尔当块的极小多项式是 \((x-\lambda)^k\)（注意这由以下事实立即推出：每个初等若尔当矩阵给出在零化子为 \((x-\lambda)^k\) 的循环 \(F[x]\)-子模上的作用）。由于 \(A\) 与 \(B\) 有相同的极小多项式，\((x-\lambda)^k\) 的最小公倍不能有重根。由此 \(k\) 必为 1，即每个若尔当块的大小都是 1，\(B\) 是对角矩阵。
\end{proof}

\noindent\textbf{从一种标准形换到另一种}

我们继续假设域 \(F\) 包含 \(T\)（或 \(A\)）的所有特征值，故在 \(F\) 上有理标准形与若尔当标准形都存在。从一种形式换到另一种的过程与第 5.2 节中关于有限阿贝尔群描述的算法完全相同（那里由不变因子表确定初等因子，反之亦然）。

简要重述：初等因子是不变因子的素幂因子。把它们从不变因子得到的方法是：把每个不变因子写成互异线性因子的幂之积；所得的线性多项式的幂的集合就是初等因子集合。例如若 \(T\) 的不变因子是
\[
(x-1)(x-3)^3, \quad (x-1)(x-2)(x-3)^3, \quad (x-1)(x-2)^2(x-3)^3,
\]
则初等因子是
\[
(x-1),\ (x-3)^3,\ (x-1),\ (x-2),\ (x-3)^3,\ (x-1),\ (x-2)^2,\ (x-3)^3.
\]
最大不变因子是初等因子中各互异素幂最大者的乘积，次大不变因子是剩余初等因子中各互异素幂最大者的乘积，依此类推。给定初等因子表，我们可以如下求不变因子表：首先把这些初等因子按特征值分成 \(n\) 个表（每个特征值一个）。在每个表中按次数递增（即非递减）排列这些多项式。然后通过补上适当个数的常数多项式 1 使所有表有相同的长度。最后把每个表中第 \(i\) 个多项式相乘得到第 \(i\) 个不变因子。例如，若 \(T\) 的初等因子是
\[
(x-1)^3,\ (x-1)^3,\ (x-1)^4,\ (x-1)^5,\ (x+4),\ (x+4)^2,\ (x+4)^3,\ (x-5)^2,\ (x-5)^3,
\]
则中间的三个表是
\[
(x-1)^3,\ (x-1)^3,\ (x-1)^4,\ (x-1)^5;
\]
\[
1,\ x+4,\ (x+4)^2,\ (x+4)^3;
\]
\[
1,\ 1,\ (x-5)^2,\ (x-5)^3,
\]
故不变因子表是
\[
(x-1)^3,\quad (x-1)^3(x+4),\quad (x-1)^4(x+4)^2(x-5)^2,\quad (x-1)^5(x+4)^3(x-5)^3.
\]

\noindent\textbf{初等因子分解算法：化为若尔当标准形}

定理 \ref{thm:12.21} 指出了确定任何给定矩阵 \(A\) 的不变因子的计算步骤。把这些不变因子分解就得到 \(A\) 的初等因子，从而如上确定 \(A\) 的若尔当标准形。

定理 \ref{thm:12.21} 之后的不变因子分解算法从 \(V\) 的一组基 \(e_1, \ldots, e_n\) 出发，产生 \(V\) 的一组元素 \(f_1, \ldots, f_m\)，它们是 \(V\) 的不变因子分解中各循环因子的 \(F[x]\)-模生成元（零化子分别为 \((a_1(x)), \ldots, (a_m(x))\)）。由于初等因子分解是通过把中国剩余定理应用于循环模 \(F[x]/(a_i(x))\) 而从不变因子分解得到的，这给出 \(V\) 的初等因子分解中各循环因子的一组 \(F[x]\)-模生成元。这些元素随后给出 \(V\) 的一组显式向量空间基，关于它对应于 \(A\) 的线性变换处于若尔当标准形（等价地，给出显式矩阵 \(P\) 使 \(P^{-1}AP\) 处于若尔当标准形）。与不变因子分解算法一样，我们首先在分解向量空间的一般背景下陈述结果，然后描述把给定 \(n \times n\) 矩阵 \(A\) 化为若尔当标准形的算法。

\noindent\textbf{初等因子分解算法}

（1）到（3）：算法的前三步是定理 \ref{thm:12.21} 之后的不变因子分解算法中的步骤。

（4）对为 \(A\) 计算出的每个不变因子 \(a(x)\)，写成
\[
a(x) = (x-\lambda_1)^{\alpha_1}(x-\lambda_2)^{\alpha_2}\cdots(x-\lambda_s)^{\alpha_s},
\]
其中 \(\lambda_1, \ldots, \lambda_s \in F\) 互异。设 \(f \in V\) 是（3）中计算出的对应于不变因子 \(a(x)\) 的循环因子的 \(F[x]\)-模生成元。则元素
\[
\frac{a(x)}{(x-\lambda_1)^{\alpha_1}}f,\quad \ldots,\quad \frac{a(x)}{(x-\lambda_s)^{\alpha_s}}f
\]
（注意 \(\frac{a(x)}{(x-\lambda_i)^{\alpha_i}} \in F[x]\) 是多项式）是 \(V\) 对应于初等因子 \((x-\lambda_1)^{\alpha_1}, \ldots, (x-\lambda_s)^{\alpha_s}\) 的各循环因子的 \(F[x]\)-模生成元。

（5）若 \(g_i = \frac{a(x)}{(x-\lambda_i)^{\alpha_i}}f\) 是对应于初等因子 \((x-\lambda_i)^{\alpha_i}\) 的 \(V\) 的循环因子的 \(F[x]\)-模生成元，则这个循环因子相应的向量空间基由元素
\[
(T-\lambda_i)^{\alpha_i-1}g_i, \quad \ldots, \quad (T-\lambda_i)g_i, \quad g_i
\]
给出。

（6）把（5）中计算出的向量空间基的第 \(k\) 个元素用原来的向量空间基 \([e_1, e_2, \ldots, e_n]\) 表示，并把这些坐标用作 \(n \times n\) 矩阵 \(P\) 的第 \(k\) 列。则 \(P^{-1}AP\) 处于若尔当标准形（若尔当块按（5）中对 \(V\) 的各循环因子所用的次序排列）。

\noindent\textbf{把 \(n \times n\) 矩阵化为若尔当标准形}

（1）到（2）：前两步是定理 \ref{thm:12.21} 之后把 \(n \times n\) 矩阵化为有理标准形算法的步骤。

（3）当 \(xI-A\) 已对角化为定理 \ref{thm:12.21} 中的形式时，矩阵 \(P'\) 的前 \(n-m\) 列为 0（这提供有用的数值检验），其余 \(m\) 列非零。对每个依次的 \(i = 1, 2, \ldots, m\)：

（a）把第 \(i\) 个非恒定对角元素（其次数为 \(d_i\)）分解为
\[
a(x) = (x-\lambda_1)^{\alpha_1}(x-\lambda_2)^{\alpha_2}\cdots(x-\lambda_s)^{\alpha_s},
\]
其中 \(\lambda_1, \ldots, \lambda_s \in F\) 互异（这里 \(a(x) = a_i(x)\) 是第 \(i\) 个非恒定对角元素，\(s\) 依赖于 \(i\)）。

（b）把 \(P'\) 的第 \(i\) 个非零列依次乘以 \(d_i\) 个矩阵
\[
(A-\lambda_1I)^{\alpha_1-1}(A-\lambda_2I)^{\alpha_2}\cdots(A-\lambda_sI)^{\alpha_s},
\]
\[
(A-\lambda_1I)^{\alpha_1-1}(A-\lambda_2I)^{\alpha_2}\cdots(A-\lambda_sI)^{\alpha_s}, \quad \ldots, \quad (A-\lambda_1I)^{\alpha_1}(A-\lambda_2I)^{\alpha_2}\cdots(A-\lambda_sI)^{\alpha_s-1}, \quad (A-\lambda_1I)^{\alpha_1}\cdots(A-\lambda_sI)^{\alpha_s}.
\]
（按原文，这些因子依次把各 \(\alpha_i\) 逐项减 1，最后一个是 \(\prod_i(A-\lambda_iI)^{\alpha_i}\)，其中（按（5）的次序）这些矩阵依次为 \(\prod_{j\neq s}(A-\lambda_jI)^{\alpha_j}\cdot(A-\lambda_sI)^{\alpha_s-1}, \ldots, \prod_{j\neq1}(A-\lambda_jI)^{\alpha_j}\cdot(A-\lambda_1I)^{\alpha_1-1}, \prod_j(A-\lambda_jI)^{\alpha_j}\) 的各种排列，见正文中列出的 \(d_i\) 个矩阵。）

（c）按此顺序把（b）中所得的列向量用作 \(n \times n\) 矩阵 \(P\) 接下来的 \(d_i\) 列。

则 \(P^{-1}AP\) 处于若尔当标准形（其若尔当块对应于（a）中因子的次序）。

\begin{example}
（1）设 \(A, B, C\) 是上一节例 1 中的矩阵，\(F = \mathbb{Q}\). 注意 \(\mathbb{Q}\) 包含这些矩阵的所有特征值。由于我们已经确定了这些矩阵的不变因子，可以立即得到它们的初等因子。\(A\) 的初等因子是 \(x-2, x-2\) 与 \(x-3\)，而 \(B\) 与 \(C\) 的初等因子是 \((x-2)^2\) 与 \(x-3\)，故相应的若尔当标准形分别是
\[
\begin{pmatrix} 2 & 0 & 0 \\ 0 & 2 & 0 \\ 0 & 0 & 3 \end{pmatrix}, \qquad \begin{pmatrix} 2 & 1 & 0 \\ 0 & 2 & 0 \\ 0 & 0 & 3 \end{pmatrix}, \qquad \begin{pmatrix} 2 & 1 & 0 \\ 0 & 2 & 0 \\ 0 & 0 & 3 \end{pmatrix}.
\]
注意 \(A\) 相似于对角矩阵，但由推论 \ref{cor:12.25}，\(B\) 与 \(C\) 不能对角化。

（2）对矩阵 \(A\)，我们在上一节例 2 中确定 \(f_1 = -7e_1+7e_2+e_3\) 与 \(f_2 = -e_1+e_2\) 是 \(V\) 在其不变因子分解中两个循环因子的 \(\mathbb{Q}[x]\)-模生成元，分别对应于不变因子 \(x-2\) 与 \((x-2)(x-3)\). 用上面描述的第一个算法，元素 \(f_1\)、\((x-3)f_2\) 与 \((x-2)f_2\) 因此是 \(V\) 在其初等因子分解中三个循环因子的 \(\mathbb{Q}[x]\)-模生成元，分别对应于初等因子 \(x-2\)、\(x-2\) 与 \(x-3\). 简单计算表明这些元素分别是 \(-7e_1+7e_2+e_3\)、\(-e_1\) 与 \(-2e_1+e_2\). 于是矩阵
\[
P = \begin{pmatrix} -7 & -1 & -2 \\ 7 & 0 & 1 \\ 1 & 0 & 0 \end{pmatrix}
\]
把 \(A\) 共轭为它的若尔当标准形
\[
P^{-1}AP = \begin{pmatrix} 2 & 0 & 0 \\ 0 & 2 & 0 \\ 0 & 0 & 3 \end{pmatrix},
\]
容易验证。这个矩阵的各列也可以按上面第二个算法用上一节例 2 中计算出的矩阵 \(P'\) 的非零列得到，结果再次给出同一个矩阵 \(P\).

（3）对上一节例 3 中的 \(4 \times 4\) 矩阵 \(D\)，不变因子是 \((x-1)^2, (x-1)^2\)，相应的 \(\mathbb{Q}[x]\)-模生成元分别是 \(f_1 = e_1\) 与 \(f_2 = e_2\). 这些同时也是这个矩阵的初等因子。这两个因子相应的向量空间基分别由 \((T-I)f_1, f_1\) 与 \((T-I)f_2, f_2\) 给出。简单计算表明这些元素分别是 \(2e_2+e_3, e_1\) 与 \(2e_1-e_2+e_4, e_2\). 于是矩阵
\[
P = \begin{pmatrix} 0 & 1 & 2 & 0 \\ 2 & 0 & -2 & 1 \\ 1 & 0 & 0 & 0 \\ 0 & 0 & 1 & 0 \end{pmatrix}
\]
把 \(D\) 共轭为它的若尔当标准形
\[
P^{-1}DP = \begin{pmatrix} 1 & 1 & 0 & 0 \\ 0 & 1 & 0 & 0 \\ 0 & 0 & 1 & 1 \\ 0 & 0 & 0 & 1 \end{pmatrix},
\]
容易验证。这个矩阵的各列也可以按上面第二个算法用上一节例 3 中计算出的 \(P'\) 的非零列得到，结果再次给出同一个矩阵 \(P\).

（4）元素取自 \(\mathbb{C}\)、特征多项式为 \((x^4-1)(x^2-1)\) 的 \(6 \times 6\) 矩阵的相似类集合由上一组例子中有理标准形所代表的 4 个类组成（在 \(\mathbb{C}\) 上没有额外的不变因子表）。然而它们的若尔当标准形不能都写在 \(\mathbb{Q}\) 上。例如若不变因子是 \((x-1)(x+1)\) 与 \((x-1)(x+1)(x^2+1)\)，则初等因子是 \(x-1, x+1, x-1, x+1, x-i, x+i\)，其中 \(i\) 是 \(\mathbb{C}\) 中 \(-1\) 的一个平方根，故这个矩阵的若尔当形是对角元素为 \(1, 1, -1, -1, i, -i\) 的对角矩阵。

（5）相比之下，\(\mathbb{C}\) 上满足 \(A^6 = I\) 的 \(3 \times 3\) 矩阵 \(A\) 的相似类集合比 \(\mathbb{Q}\) 上的大得多。若 \(A\) 是这样的矩阵，则 \(m_A(x) \mid x^6-1\)，而由于后者在 \(\mathbb{C}\) 中没有重根，\(A\) 的极小多项式没有重根。由推论 \ref{cor:12.25}，\(A\) 的若尔当标准形是对角矩阵。由于这个对角矩阵有相同的极小多项式，它的 6 次幂也是恒等矩阵，故每个对角元素都是 6 次单位根。对每个由 6 次单位根组成的三元组 \(\{\zeta_1, \zeta_2, \zeta_3\}\) 我们得到一个若尔当标准形，而两个这样的形式相同（即给出相似的矩阵）当且仅当这些三元组互为置换。可以算出在相似意义下这样的 \(A\) 共有 56 个相似类。
\end{example}

\subsection*{习题}

\begin{enumerate}
\item 假设向量空间 \(V\) 是零化子分别为 \((x+1)^2\)、\((x-1)(x^2+1)^2\)、\((x^4-1)\) 与 \((x+1)(x^2-1)\) 的循环 \(F[x]\)-模的直和。确定 \(V\) 的不变因子与初等因子。

\item 证明若 \(\lambda_1, \ldots, \lambda_n\) 是 \(n \times n\) 矩阵 \(A\) 的特征值，则对任意 \(k \geq 0\)，\(\lambda_1^k, \ldots, \lambda_n^k\) 是 \(A^k\) 的特征值。

\item 用上面例 2 的方法确定显式矩阵 \(P_1\) 与 \(P_2\)，使 \(P_1^{-1}BP_1\) 与 \(P_2^{-1}CP_2\) 处于若尔当标准形。用它显式构造一个把 \(B\) 共轭为 \(C\) 的矩阵 \(Q\)（直接证明这两个矩阵相似）。

\item 证明矩阵
\[
\begin{pmatrix} 9 & 4 & 5 \\ -4 & 0 & -3 \\ -6 & -4 & -2 \end{pmatrix}
\]
的若尔当标准形就是本章开头所陈述的矩阵。显式确定把这个矩阵共轭为其若尔当标准形的矩阵 \(P\). 解释为什么这个矩阵不能对角化。

\item 计算矩阵
\[
\begin{pmatrix} 1 & 1 & -1 \\ 0 & 1 & 3 \\ 0 & 0 & 3 \end{pmatrix}
\]
的若尔当标准形。

\item 判断下列矩阵中哪些相似：
\[
\begin{pmatrix} 4 & 4 & 2 \\ 1 & -3 & 2 \\ -1 & 3 & -4 \end{pmatrix}.
\]

\item 确定下列矩阵的若尔当标准形：
\[
\begin{pmatrix} -1 & 1 & 0 \\ -4 & 3 & 0 \\ 1 & -1 & -3 \end{pmatrix}, \qquad \begin{pmatrix} -3 & -1 & -3 \\ -4 & -4 & -4 \\ 0 & 0 & 0 \end{pmatrix}.
\]

\item 证明矩阵
\[
A = \begin{pmatrix} 6 & 0 & 3 \\ -4 & 2 & -4 \\ -2 & 0 & 1 \end{pmatrix}, \qquad B = \begin{pmatrix} -1 & 2 & -1 \\ -10 & 6 & -14 \\ -6 & 3 & -7 \end{pmatrix}
\]
相似。证明 \(A\) 与 \(B\) 都可以对角化，并确定显式矩阵 \(P_1, P_2\) 使 \(P_1^{-1}AP_1\) 与 \(P_2^{-1}BP_2\) 处于对角形。

\item 证明矩阵
\[
A = \begin{pmatrix} 1 & 1 & -1 \\ 3 & 2 & 0 \end{pmatrix}, \qquad B = \begin{pmatrix} -1 & -3 & 2 \\ 3 & -1 & 4 \\ -2 & 5 & -2 \end{pmatrix}
\]
都以 \((x-1)^2(x+1)\) 为特征多项式，但其中一个可以对角化而另一个不能。确定这两个矩阵的若尔当标准形。（按原文，\(A, B\) 均为 \(3 \times 3\) 矩阵：\(A = \begin{pmatrix}1&1&-1\\-3&2&-4\\3&2&0\end{pmatrix}\) 与 \(B = \begin{pmatrix}-1&-3&2\\-4&3&4\\-2&5&-2\end{pmatrix}\) 一类；请对照原书。）

\item 求 \(\mathbb{C}\) 上 \(2 \times 2\)、\(3 \times 3\) 与 \(4 \times 4\) 矩阵的全部若尔当标准形。

\item 验证矩阵
\[
A = \begin{pmatrix} 1 & 0 & 0 & 0 \\ 0 & 1 & 0 & 0 \\ -2 & -2 & 0 & 1 \\ -2 & 0 & -1 & -1 \end{pmatrix}
\]
的特征多项式在 \(\mathbb{Q}\) 上是线性因子之积。确定 \(A\) 在 \(\mathbb{Q}\) 上的有理标准形与若尔当标准形。

\item 确定矩阵
\[
\begin{pmatrix} 1 & 2 & 0 & 0 \\ 0 & 1 & 2 & 0 \\ 0 & 0 & 1 & 2 \\ 0 & 0 & 0 & 1 \end{pmatrix}
\]
的若尔当标准形。

\item 确定矩阵
\[
\begin{pmatrix} 3 & 0 & -2 & -3 \\ 4 & -8 & 14 & -15 \\ 2 & -4 & 7 & -7 \\ 0 & 2 & -4 & 3 \end{pmatrix}
\]
的若尔当标准形。

\item 证明矩阵
\[
A = \begin{pmatrix} 0 & 0 & -4 & -7 \\ 1 & -8 & 15 & -13 \\ -1 & -4 & 7 & -7 \\ -2 & 4 & 2 & -5 \end{pmatrix}
\]
相似。（按原文，本题给出两个相似矩阵，其第二个在扫描件中无法辨认，请对照原书。）

\item 证明矩阵
\[
A = \begin{pmatrix} 5 & 2 & -8 & -8 \\ -6 & -3 & 8 & 8 \\ -3 & -1 & 3 & 4 \\ 3 & 1 & -4 & -5 \end{pmatrix}, \qquad B = \begin{pmatrix} 1 & 0 & 0 & 0 \\ 0 & 1 & 0 & 0 \\ 0 & 0 & 3 & 0 \\ 0 & 0 & 0 & 3 \end{pmatrix}
\]
都以 \((x-3)(x+1)^3\) 为特征多项式。判断它们是否相似，并确定每个矩阵的若尔当标准形。

\item 确定矩阵
\[
\begin{pmatrix} 0 & 0 & 1 & 1 \\ 0 & 0 & 0 & 1 \\ 0 & 0 & 0 & 0 \\ 0 & 0 & 0 & 0 \end{pmatrix}
\]
的若尔当标准形，并确定把它共轭为其若尔当标准形的矩阵 \(P\).

\item 证明任何矩阵 \(A\) 相似于它的转置 \(A^t\).

\item 确定特征多项式为 \((x-2)^3(x-3)^2\) 的线性变换的全部可能的若尔当标准形。

\item 证明所有具有特征多项式 \(f(x)\) 的 \(n \times n\) 矩阵都相似，当且仅当 \(f(x)\) 在 \(F[x]\) 中的唯一分解没有重复因子。

\item 证明下列矩阵在 \(M_p(\mathbb{F}_p)\)（元素取自 \(\mathbb{F}_p\) 的 \(p \times p\) 矩阵）中相似：
\[
\begin{pmatrix} 0 & 1 & 0 & 0 & 0 \\ 0 & 0 & 1 & 0 & 0 \\ 0 & 0 & 0 & 1 & 0 \\ 0 & 0 & 0 & 0 & 1 \\ 1 & 1 & 0 & 0 & 0 \end{pmatrix}, \qquad \begin{pmatrix} 0 & 1 & 0 & 0 & 0 \\ 0 & 0 & 1 & 0 & 0 \\ 0 & 0 & 0 & 1 & 0 \\ 0 & 0 & 0 & 0 & 1 \\ 0 & 0 & 1 & 0 & 0 \end{pmatrix}.
\]

\item 证明若 \(A^2 = A\)，则 \(A\) 相似于对角线上只有 0 与 1 的对角矩阵。

\item 证明元素取自 \(\mathbb{C}\)、满足 \(A^3 = A\) 的 \(n \times n\) 矩阵 \(A\) 可以对角化。对任何域 \(F\) 同样的断言成立吗？

\item 假设 \(A\) 是元素取自 \(\mathbb{Q}\)、满足 \(A^3 = I\) 但 \(A \neq I\) 的 \(2 \times 2\) 矩阵。把 \(A\) 写成有理标准形，并把它看作 \(\mathbb{C}\) 上的矩阵写成若尔当标准形。

\item 证明不存在满足 \(A^8 = I\) 但 \(A \neq I\) 的 \(\mathbb{Q}\) 上 \(3 \times 3\) 矩阵 \(A\).

\item 确定元素全为 1 的 \(\mathbb{Q}\) 上 \(n \times n\) 矩阵的若尔当标准形。

\item 确定元素全为 1 的 \(\mathbb{F}_p\) 上 \(n \times n\) 矩阵的若尔当标准形（答案取决于 \(p\) 是否整除 \(n\)）。

\item 确定除主对角线元素全为 0 之外其余元素全为 1 的 \(\mathbb{Q}\) 上 \(n \times n\) 矩阵的若尔当标准形。

\item 确定除主对角线元素全为 0 之外其余元素全为 1 的 \(\mathbb{F}_p\) 上 \(n \times n\) 矩阵的若尔当标准形。
\end{enumerate}

\noindent 对应于 \(V\) 的所有作为同一个 \(x-\lambda\) 的幂的初等因子的循环子模的直和称为 \(T\) 对应于特征值 \(\lambda\) 的\textbf{广义特征空间}。注意这是 \(V\) 关于 \(F[x]\) 的素元 \(p = x-\lambda\) 的 \(p\)-准素分支，由被线性变换 \(T-\lambda\) 的某个幂零化的 \(V\) 中元素组成。\(T\) 在 \(\lambda\) 的广义特征空间上的矩阵是 \(T\) 的所有特征值为 \(\lambda\) 的若尔当块组成的分块对角矩阵。

\begin{enumerate}
\setcounter{enumi}{28}
\item 假设 \(V_i\) 是 \(T\) 对应于特征值 \(\lambda_i\) 的广义特征空间。证明对任意 \(k \geq 0\)，\(T-\lambda_i\) 在子空间 \((T-\lambda_i)^kV_i\) 上的零化度等于 \(T-\lambda_i\) 在 \((T-\lambda_i)^kV_i\) 上的零化度，且等于 \(T\) 的特征值为 \(\lambda_i\)、大小大于 \(k\) 的若尔当块的个数（故 \(k = 0\) 时这给出若尔当块的个数）。

\item 设 \(\lambda\) 是 \(T\) 在域 \(F\) 上有限维向量空间 \(V\) 上的特征值。设 \(r_k = \dim_F (T-\lambda)^kV\) 是线性变换 \((T-\lambda)^k\) 在 \(V\) 上的秩。证明对任意 \(k \geq 1\)，\(r_{k-1}-2r_k+r_{k+1}\) 是 \(T\) 对应于 \(\lambda\)、大小为 \(k\) 的若尔当块的个数〔利用第 12.1 节习题 12〕。（这给出通过计算表示 \(T\) 的矩阵 \(A\) 的矩阵 \((A-\lambda I)^k\) 的秩来确定 \(T\) 的若尔当标准形的有效方法，参见第 11.2 节习题 31（a）。）

\item 设 \(N\) 是系数在域 \(F\) 中的 \(n \times n\) 矩阵。若 \(N\) 的某个幂是零矩阵，即对某个 \(k\) 有 \(N^k = 0\)，则称 \(N\) 是\textbf{幂零的}。证明任何幂零矩阵相似于一个分块对角矩阵，其各块是第一条上对角线为 1、其余位置为 0 的矩阵。

\item 证明若 \(N\) 是 \(n \times n\) 幂零矩阵，则事实上 \(N^n = 0\).

\item 设 \(A\) 是严格上三角的 \(n \times n\) 矩阵（主对角线上及其以下的元素全为零）。证明 \(A\) 是幂零的。

\item 证明幂零 \(n \times n\) 矩阵的迹为 0（回忆矩阵的迹是对角元素之和）。

\item 对 \(0 \leq i \leq n\)，设 \(d_i\) 是 \(xI-A\) 的所有 \(i \times i\) 子式的行列式的最大公因子（其中 \(A\) 如定理 \ref{thm:12.21} 中那样，并取 \(0 \times 0\) 子式为 1）。证明 \(A\) 的 Smith 标准形对角线上第 \(i\) 个元素是 \(d_i/d_{i-1}\). 这给出 \(A\) 的不变因子。〔证明这些最大公因子在初等行运算与列运算下不变。〕

\item 设 \(V = \mathbb{C}^n\) 是复数 \(n\) 元组 \((a_1, a_2, \ldots, a_n)\) 的通常 \(n\) 维向量空间。设 \(T\) 是由令 \(T(a_1, a_2, \ldots, a_n)\) 等于 \((0, a_1, a_2, \ldots, a_{n-1})\) 定义的线性变换。确定 \(T\) 的若尔当标准形。

\item 设 \(J\) 是 \(\mathbb{C}\) 上特征值为 \(\lambda\) 的大小为 \(n\) 的若尔当块。

（a）证明若 \(\lambda \neq 0\)，则矩阵 \(J^2\) 的若尔当标准形是特征值为 \(\lambda^2\) 的大小为 \(n\) 的若尔当块。

（b）若 \(\lambda = 0\)，证明 \(J^2\) 的若尔当标准形有两个块（特征值为 0），当 \(n\) 为偶数时大小均为 \(n/2\)；当 \(n\) 为奇数时大小分别为 \((n-1)/2\) 与 \((n+1)/2\).

\item 确定 \(A \in M_n(\mathbb{C})\) 有平方根（即存在另一个矩阵 \(B \in M_n(\mathbb{C})\) 使 \(A = B^2\)）的充分必要条件。〔假设 \(B\) 处于若尔当标准形，并利用上一习题考虑 \(B^2\) 的若尔当标准形。〕

\item 设 \(J\) 是特征为 2 的域 \(F\) 上特征值为 \(\lambda\) 的大小为 \(n\) 的若尔当块。确定矩阵 \(J^2\) 的若尔当标准形。确定 \(A \in M_n(F)\) 有平方根（即存在另一个矩阵 \(B \in M_n(F)\) 使 \(A = B^2\)）的充分必要条件。
\end{enumerate}

\noindent 余下的习题探讨矩阵的函数（幂级数），并介绍若尔当标准形在微分方程理论中的一些应用。这些习题中矩阵都假设是元素取自域 \(K\) 的 \(n \times n\) 矩阵，其中 \(K\) 是实数域或复数域。设 \(G(x) = \sum_{k=0}^{\infty}a_kx^k\) 是系数在 \(K\) 中的幂级数。设 \(G_N(x) = \sum_{k=0}^{N}a_kx^k\) 是 \(G(x)\) 的第 \(N\) 个部分和，对每个 \(A \in M_n(K)\) 设 \(G_N(A)\) 是在这个多项式中代入 \(A\) 所得的 \(M_n(K)\) 的元素。对每个固定的 \(i, j\)，令 \(c_{ij}^N\) 为矩阵 \(G_N(A)\) 的 \(i, j\) 元素，得到实数或复数序列 \(c_{ij}^N\)（\(N = 0, 1, 2, \ldots\)）。若对每个 \(i, j \in \{1, 2, \ldots, n\}\)，序列 \(c_{ij}^N\) 收敛到 \(C\) 的 \(i, j\) 元素，则称级数 \(G(A) = \sum_{k=0}^{\infty}a_kA^k\) 收敛到矩阵 \(C \in M_n(K)\)（此时记 \(G(A) = C\)）。若存在 \(C \in M_n(K)\) 使 \(G(A) = C\)，则说 \(G(A)\) 收敛。若 \(A\) 是 \(1 \times 1\) 矩阵，这就是 \(K\) 中级数收敛的通常概念。

\noindent 对 \(A = (a_{ij}) \in M_n(K)\) 定义 \(\|A\| = \sum_{i,j=1}^n |a_{ij}|\)，即 \(\|A\|\) 是 \(A\) 的所有元素的绝对值之和。

\begin{enumerate}
\setcounter{enumi}{39}
\item 证明对所有 \(A, B \in M_n(K)\) 与所有 \(a \in K\)：（a）\(\|A+B\| \leq \|A\|+\|B\|\)；（b）\(\|AB\| \leq \|A\| \cdot \|B\|\)；（c）\(\|aA\| = |a| \cdot \|A\|\).

\item 设 \(R\) 是实或复幂级数 \(G(x)\) 的收敛半径（若 \(G(x)\) 对所有 \(x \in K\) 收敛则 \(R = \infty\)）。

（a）证明若 \(\|A\| < R\)，则 \(G(A)\) 收敛。

（b）由此断定对所有矩阵 \(A\)，下列幂级数收敛：
\[
\sin(A) = A-\frac{A^3}{3!}+\frac{A^5}{5!}+\cdots+(-1)^k\frac{A^{2k+1}}{(2k+1)!}+\cdots,
\]
\[
\cos(A) = I-\frac{A^2}{2!}+\frac{A^4}{4!}+\cdots+(-1)^k\frac{A^{2k}}{(2k)!}+\cdots,
\]
\[
\exp(A) = I+A+\frac{A^2}{2!}+\frac{A^4}{4!}+\cdots+\frac{A^k}{k!}+\cdots,
\]
其中 \(I\) 是 \(n \times n\) 恒等矩阵。
\end{enumerate}

\noindent 鉴于在微分方程理论中的应用，此时我们引入变量 \(t\)，使得对 \(A \in M_n(K)\)，矩阵 \(At\) 由把 \(A\) 的每个元素乘以 \(t\) 得到（这与把 \(A\) 乘以“纯量”矩阵 \(tI\) 相同）。我们得到由 \(t \mapsto G(At)\) 定义的从 \(K\) 的子集到 \(M_n(K)\) 的函数，定义在所有使级数 \(G(At)\) 收敛的点 \(t\) 处。特别地，\(\sin(At)\)、\(\cos(At)\) 与 \(\exp(At)\) 对所有 \(t \in K\) 收敛。

\begin{enumerate}
\setcounter{enumi}{41}
\item 设 \(P\) 是非奇异 \(n \times n\) 矩阵。

（a）证明 \(PG(At)P^{-1} = G(PAtP^{-1}) = G(PAP^{-1}t)\).（这意味着在基变换的意义下，只需对标准形中的矩阵 \(A\) 计算 \(G(At)\)。）〔对部分和取极限即得第一个等式。第二个等式由矩阵 \(tI\) 与所有矩阵交换立即推出。〕

（b）证明若 \(A\) 是矩阵 \(A_1, A_2, \ldots, A_m\) 的直和，则 \(G(At)\) 是矩阵 \(G(A_1t), G(A_2t), \ldots, G(A_mt)\) 的直和。

（c）证明若 \(Z\) 是对角元素为 \(z_1, z_2, \ldots, z_n\) 的对角矩阵，则 \(G(Zt)\) 是对角元素为 \(G(z_1t), G(z_2t), \ldots, G(z_nt)\) 的对角矩阵。
\end{enumerate}

\noindent 习题 41（b）中定义的矩阵 \(\exp(A)\) 称为 \(A\) 的\textbf{指数}，常记作 \(e^A\). 下面三个习题导出初等若尔当矩阵 \(J\) 的矩阵 \(\exp(Jt)\) 的公式。

\begin{enumerate}
\setcounter{enumi}{42}
\item 证明若 \(A\) 与 \(B\) 是可交换矩阵，则 \(\exp(A+B) = \exp(A)\exp(B)\).〔把 \(A\) 与 \(B\) 看作可交换的未定元，通过比较左端的幂级数与右端两个幂级数之积得出。〕

\item 用上一习题证明：若 \(M\) 是任意矩阵、\(\lambda\) 是 \(K\) 的任意元素，则 \(\exp(\lambda tI+M) = e^{\lambda t}\exp(M)\).

\item 设 \(N\) 是第一上对角线为 1、其余位置为 0 的 \(r \times r\) 矩阵。计算下列幂零 \(r \times r\) 矩阵的指数：若 \(Nt\)，则 \(\exp(Nt)\) 是主对角线与第一条上对角线上的元素依次为 \(1, t, t^2/2!, \ldots, t^{r-1}/(r-1)!\) 的上三角矩阵（其余位置为 0）。由此断定若 \(J\) 是特征值为 \(\lambda\) 的 \(r \times r\) 初等若尔当矩阵，则 \(\exp(Jt)\) 是上三角矩阵，其主对角线元素全为 \(e^{\lambda t}\)、第一条上对角线元素依次为 \(te^{\lambda t}, \frac{t^2}{2!}e^{\lambda t}, \ldots, \frac{t^{r-1}}{(r-1)!}e^{\lambda t}\).〔第一部分利用 \(Nt\) 是幂零矩阵这一观察：由于 \(Nt\) 幂零，\(\exp(Nt)\) 是 \(Nt\) 的多项式，即幂级数中除有限多项外全为 0. 为计算 \(Jt\) 的指数，把 \(Jt\) 写成 \(\lambda tI+Nt\) 并用习题 44 取 \(M = Nt\).〕
\end{enumerate}

\noindent 设 \(A \in M_n(K)\)，\(P\) 是使 \(P^{-1}AP\) 处于若尔当标准形的基变换矩阵。假设 \(P^{-1}AP\) 是初等若尔当矩阵 \(J_1, \ldots, J_m\) 之和。上述习题（取 \(t = 1\)）表明 \(\exp(A)\) 可以容易地求出：把 \(E = \exp(P^{-1}AP)\) 写成矩阵 \(\exp(J_1), \ldots, \exp(J_m)\) 的直和，然后换回基得到 \(\exp(A) = PEP^{-1}\).

\begin{enumerate}
\setcounter{enumi}{45}
\item 对本节例 3 中给出的 \(4 \times 4\) 矩阵 \(D\) 与 \(P\)，验证 \(e^D\) 与 \(\exp(D)\) 的公式（用 \(e = e^1\) 表示；按原文给出 \(\exp(D)\) 的显式矩阵）。

\item 计算下列矩阵的指数：（a）本节例 2 中的矩阵 \(A\)；（b）习题 4 中的矩阵（你已算出它的若尔当标准形与基变换矩阵）；（c）习题 16 中的矩阵。

\item 证明 \(\exp(0) = I\)（这里 0 是零矩阵、\(I\) 是恒等矩阵）。由此断定 \(\exp(A)\) 非奇异，且对所有 \(A \in M_n(K)\) 其逆为 \(\exp(-A)\).

\item 证明 \(\det(\exp(A)) = e^{\operatorname{tr}(A)}\)，其中 \(\operatorname{tr}(A)\) 是 \(A\) 的迹（对角元素之和）。

\item 固定 \(A \in M_n(K)\). 证明由 \(t \mapsto \exp(At)\) 定义的映射是群同态（这里 \(K\) 是域 \(K\) 的加法群）。（注意这推广了熟悉的从 \(K\) 到 \(K^\times\) 的指数映射，后者是 \(n = 1\) 的情形。子群 \(\{\exp(At) \mid t \in K\}\) 称为 \(GL_n(K)\) 的\textbf{单参数子群}。这些子群与指数映射在 Lie 群理论中起重要作用——\(GL_n(K)\) 是 Lie 群的一个例子。）
\end{enumerate}

\noindent 设 \(G(x)\) 是收敛半径为无限大的幂级数，固定矩阵 \(A \in M_n(K)\). 矩阵 \(G(At)\) 的元素是变量 \(t\) 的、对所有 \(t\) 都有定义的 \(K\)-值函数。设 \(c_{ij}(t)\) 是 \(G(At)\) 的 \(i, j\) 元素作为 \(t\) 的函数。\(G(At)\) 对 \(t\) 的\textbf{导数}，记作 \(\frac{d}{dt}G(At)\)，是 \(i, j\) 元素为 \(\frac{d}{dt}c_{ij}(t)\)（由对 \(G(At)\) 的每个元素求导得到）的矩阵。换句话说，若我们把 \(M_n(K)\) 与 \(K^{n^2}\) 等同（把每个 \(n \times n\) 矩阵看作 \(n^2\) 元组），则 \(t \mapsto G(At)\) 是从 \(K\) 到 \(K^{n^2}\) 的映射（即 \(t\) 的向量值函数），其导数就是这个向量值函数的通常（逐分量）导数。

\begin{enumerate}
\setcounter{enumi}{50}
\item 建立下列导数的性质：（a）若 \(G(x) = \sum_{k=0}^{\infty}a_kx^k\)，则 \(\frac{d}{dt}G(At) = A\sum_{k=1}^{\infty}ka_k(At)^{k-1}\)；（b）若 \(v\) 是元素在 \(K\) 中的 \(n \times 1\) 矩阵（常向量），则 \(\frac{d}{dt}(G(At)v) = \left(\frac{d}{dt}G(At)\right)v\).

\item 由上一习题的（a）推出 \(\frac{d}{dt}\exp(At) = A\exp(At)\).
\end{enumerate}

\noindent 现在设 \(y_1(t), \ldots, y_n(t)\) 是实变量 \(t\) 的可微函数，它们由下列带常数系数 \(a_{ij} \in K\) 的一阶线性微分方程组联系：
\[
\begin{aligned}
y_1' &= a_{11}y_1+a_{12}y_2+\cdots+a_{1n}y_n\\
y_2' &= a_{21}y_1+a_{22}y_2+\cdots+a_{2n}y_n\\
&\vdots\\
y_n' &= a_{n1}y_1+a_{n2}y_2+\cdots+a_{nn}y_n
\end{aligned}
\]
（这里撇号表示对 \(t\) 求导）。设 \(A\) 是第 \(i, j\) 元素为 \(a_{ij}\) 的矩阵，故（\*) 可以写成 \(y' = Ay\)，其中 \(y\) 是函数 \(y_1(t), \ldots, y_n(t)\) 组成的列向量。元素为 \(t\) 的函数、且各列是（\*）的独立解的 \(n \times n\) 矩阵称为（\*）的\textbf{基本矩阵}。由微分方程理论，（\*）的解向量 \(y\) 的集合构成 \(K\) 上 \(n\) 维向量空间，故基本矩阵的各列是（\*）的所有解构成的向量空间的一组基。

\begin{enumerate}
\setcounter{enumi}{52}
\item 证明 \(\exp(At)\) 是（\*）的基本矩阵。再证明若 \(C\) 是元素为 \(y_1(0), \ldots, y_n(0)\) 的 \(n \times 1\) 常向量，则 \(y(t) = \exp(At)C\) 是满足初始条件 \(y(0) = C\) 的（\*）的特解。（注意这推广了一维结果：单个微分方程 \(y' = ay\) 的解空间的 1 维基为 \(e^{at}\)，而满足初始条件 \(y(0) = c\) 的唯一解是 \(y = ce^{at}\).）〔利用前面各习题。〕

\item 证明若 \(M\) 是（\*）的基本矩阵、\(Q\) 是 \(M_n(K)\) 中的非奇异矩阵，则 \(MQ\) 也是（\*）的基本矩阵。〔\(MQ\) 的各列是 \(M\) 的各列的线性组合。〕
\end{enumerate}

\noindent 现在把前两个习题应用于求解一些具体的微分方程组，方法如下：给定（\*）中的矩阵 \(A\)，计算基变换矩阵 \(P\) 使 \(B = P^{-1}AP\) 处于若尔当标准形。则 \(\exp(At) = P\exp(Bt)P^{-1}\) 是（\*）的基本矩阵。由上一习题，\(P\exp(Bt)\) 也是（\*）的基本矩阵，而 \(\exp(Bt)\) 可以用习题 45 之后的讨论中描述的方法计算（特别地，无需为得到（\*）的基本矩阵而求矩阵 \(P\) 的逆）。例如若 \(A = D\)、\(P\) 是习题 46 中给出的矩阵，则我们看到 \(A\) 的若尔当标准形是 \(B = P^{-1}AP\)，它由两个特征值为 1 的 \(2 \times 2\) 若尔当块组成。故 \(y' = Ay\) 的一个基本矩阵是 \(P\exp(Bt)\)（按原文给出其显式形式），从而给出方程组的通解（它描述了 4 维解空间）。

\begin{enumerate}
\setcounter{enumi}{54}
\item 在（a）到（c）各部分中求（\*）的基本矩阵，其中（\*）的系数矩阵 \(A\) 已给出：（a）\(A\) 是习题 47（a）中的矩阵；（b）\(A\) 是习题 47（b）中的矩阵；（c）\(A\) 是习题 47（c）中的矩阵。

\item 考虑系数矩阵 \(A\) 为习题 46 中列出的矩阵 \(D\) 的方程组（\*），其基本矩阵已在上一个习题之前算出。求满足初始条件 \(y_i(0) = 1\)（\(i = 1, 2, 3, 4\)）的（\*）的特解。
\end{enumerate}

\noindent 接下来我们探讨（\*）的一个特殊情形。给定带常数系数的 \(n\) 阶线性微分方程
\[
y^{(n)}+a_{n-1}y^{(n-1)}+\cdots+a_1y'+a_0y = 0
\]
（其中 \(y^{(k)}\) 是 \(y\) 的第 \(k\) 阶导数，\(y^{(0)} = y\)），可以令 \(y_i = y^{(i-1)}\)（\(1 \leq i \leq n\)）把它化为一个一阶线性微分方程组（这个方程组的系数矩阵在下一习题中描述）。\(n\) 阶方程（\*\*）的解构成的 \(n\) 维向量空间的一组基随后可以从线性方程组的基本矩阵得到。具体地说，在方程组的基本矩阵的 \(n \times n\) 函数列中，第一行第 \(j\) 个元素是（\*\*）的解，故基本矩阵第一行中的 \(n\) 个函数构成（\*\*）的解的一组基。

\begin{enumerate}
\setcounter{enumi}{56}
\item 证明由（\*\*）得到的 \(n\) 个一阶方程组系数矩阵 \(A\) 是多项式 \(x^n+a_{n-1}x^{n-1}+\cdots+a_1x+a_0\) 的伴侣矩阵的转置。

\item 用上述方法求下列微分方程的解空间的一组基：（a）\(y'''-3y'+2y = 0\)；（b）\(y''''+4y'''+6y''+4y'+y = 0\).
\end{enumerate}

\noindent 形如
\[
y_1' = F_1(y_1, y_2, \ldots, y_n), \quad y_2' = F_2(y_1, y_2, \ldots, y_n), \quad \ldots, \quad y_n' = F_n(y_1, y_2, \ldots, y_n)
\]
的微分方程组（其中 \(F_1, F_2, \ldots, F_n\) 是 \(n\) 个变量的函数）称为\textbf{自治系统}，我们更简洁地记作 \(y' = F(y)\)，其中 \(F = (F_1, \ldots, F_n)\).（“自治”一词意为“不依赖于时间”，它表明变量 \(t\)（可以看作时间变量）不明显地出现在右端。）方程组（\*）是每个 \(F_i\) 都是线性函数的特殊自治系统。在许多情形下，希望在不显式求出解的情况下分析自治系统解的行为（对于给定的非线性系统，能求出显式解是不太可能的）。这一研究属于自治微分方程的“定性分析”，其初步内容常在一元系统的微积分课程中处理。\(n \times n\) 自治系统定性分析的第一步是求\textbf{平衡态}，即常值解（称为平衡态是因为它们不随 \(t\) 变化）。注意常值函数 \(y = c\)（其中 \(c\) 是元素为 \(c_1, \ldots, c_n\) 的 \(n \times 1\) 常向量）是 \(y' = F(y)\) 的解当且仅当对 \(i = 1, 2, \ldots, n\) 有 \(0 = c_i' = F_i(c_1, \ldots, c_n)\)，故平衡态通过计算 \(F\) 的零点得到（对非线性系统这可能需要数值方法）。其次，给定某个解的初值，我们希望分析这个解当 \(t \to \infty\) 时的行为，这称为解的\textbf{渐近行为}。同样可能无法显式求出解，尽管由一般微分方程理论，只要函数 \(F_i\) 可微，初值问题的解唯一。平衡态 \(y = c\) 称为\textbf{全局渐近稳定}的，若每个解都当 \(t \to \infty\) 时趋于 \(c\)，即对任何解 \(y(t)\) 与所有 \(i = 1, 2, \ldots, n\) 有 \(\lim_{t\to\infty}y_i(t) = c_i\).

在线性自治系统（\*）的情形，解构成向量空间，故唯一的常值解是零解。下一个习题给出零解全局渐近稳定的一个充分条件，并给出一个例子说明线性系统的行为如何可以用其系数矩阵的特征值来分析。非线性系统可以通过在平衡态的某邻域内考虑 \(y' = Ty\)（其中 \(T\) 是在平衡态点处求值的 \(F\) 的 \(n \times n\) Jacobi 矩阵）来用线性系统近似。这样，线性系统的分析在一般自治系统的局部分析中起重要作用。

\begin{enumerate}
\setcounter{enumi}{58}
\item 证明若 \(A\) 的所有特征值都有负实部，则由对所有 \(i \in \{1, \ldots, n\}\) 令 \(y_i(t) = 0\) 给出的（\*）的解（即零解）是全局渐近稳定的。〔对不熟悉复指数函数行为者，可假设所有特征值都是实数（从而是负实数）。利用解的显式性质证明它们都当 \(t \to \infty\) 时趋于零。〕
\end{enumerate}
